% concatenated from sn-article.tex
%Version 3.1 December 2024
% See section 11 of the User Manual for version history
%
%%%%%%%%%%%%%%%%%%%%%%%%%%%%%%%%%%%%%%%%%%%%%%%%%%%%%%%%%%%%%%%%%%%%%%
%%                                                                 %%
%% Please do not use \input{...} to include other tex files.       %%
%% Submit your LaTeX manuscript as one .tex document.              %%
%%                                                                 %%
%% All additional figures and files should be attached             %%
%% separately and not embedded in the \TeX\ document itself.       %%
%%                                                                 %%
%%%%%%%%%%%%%%%%%%%%%%%%%%%%%%%%%%%%%%%%%%%%%%%%%%%%%%%%%%%%%%%%%%%%%

%%\documentclass[referee,sn-basic]{sn-jnl}% referee option is meant for double line spacing

%%=======================================================%%
%% to print line numbers in the margin use lineno option %%
%%=======================================================%%

%%\documentclass[lineno,pdflatex,sn-basic]{sn-jnl}% Basic Springer Nature Reference Style/Chemistry Reference Style

%%=========================================================================================%%
%% the documentclass is set to pdflatex as default. You can delete it if not appropriate.  %%
%%=========================================================================================%%

%%\documentclass[sn-basic]{sn-jnl}% Basic Springer Nature Reference Style/Chemistry Reference Style

%%Note: the following reference styles support Namedate and Numbered referencing. By default the style follows the most common style. To switch between the options you can add or remove “Numbered” in the optional parenthesis. 
%%The option is available for: sn-basic.bst, sn-chicago.bst%  
 
%%\documentclass[pdflatex,sn-nature]{sn-jnl}% Style for submissions to Nature Portfolio journals
%%\documentclass[pdflatex,sn-basic]{sn-jnl}% Basic Springer Nature Reference Style/Chemistry Reference Style
%%%%%%%
\documentclass[pdflatex,sn-mathphys-ay]{sn-jnl}
%%%%%
%\documentclass[pdflatex,sn-mathphys-num]{sn-jnl}% Math and Physical Sciences Numbered Reference Style
%%\documentclass[pdflatex,sn-mathphys-ay]{sn-jnl}% Math and Physical Sciences Author Year Reference Style
%%\documentclass[pdflatex,sn-aps]{sn-jnl}% American Physical Society (APS) Reference Style
%%\documentclass[pdflatex,sn-vancouver-num]{sn-jnl}% Vancouver Numbered Reference Style
%%\documentclass[pdflatex,sn-vancouver-ay]{sn-jnl}% Vancouver Author Year Reference Style
%%\documentclass[pdflatex,sn-apa]{sn-jnl}% APA Reference Style
%%\documentclass[pdflatex,sn-chicago]{sn-jnl}% Chicago-based Humanities Reference Style

%%%% Standard Packages
%%<additional latex packages if required can be included here>

\usepackage{graphicx}%
\usepackage{multirow}%
\usepackage{amsmath,amssymb,amsfonts}%
\usepackage{amsthm}%
\usepackage{mathrsfs}%
\usepackage[title]{appendix}%
\usepackage{xcolor}%
\usepackage{textcomp}%
\usepackage{manyfoot}%
\usepackage{booktabs}%
\usepackage{algorithm}%
\usepackage{algorithmicx}%
\usepackage{algpseudocode}%
\usepackage{listings}%

\usepackage{amsmath,amsfonts,amsthm,bm} % Math packages

%\usepackage{natbib}
%\usepackage[colorlinks,citecolor=blue,urlcolor=blue,filecolor=blue,backref=page]{hyperref}
\usepackage{graphicx} 
\usepackage{array}
%\usepackage[caption=false,font=normalsize,labelfont=sf,textfont=sf]{subfig}
\usepackage{textcomp}
\usepackage{stfloats}
\usepackage{mathabx}
\usepackage{url}
\usepackage{verbatim}
\usepackage{multirow}
\usepackage{multicol}
\usepackage{xcolor}
\usepackage{microtype}
\usepackage{longtable}
\usepackage{tabto}
\usepackage{tabularx}
\usepackage{amsmath}
\usepackage{float}
\usepackage{subcaption}
\usepackage{caption}
\usepackage{tikz}
\usepackage{threeparttable} 
\usepackage{todonotes}
\usepackage{placeins}

%%
\usepackage{color}
\newcommand{\lcolor}[1]{\textcolor{magenta}{#1}}  % Larry's colored text
\newcommand{\lcomm}[1]{\marginpar{\tiny\lcolor{#1}}}  % Larry's marginpars

\newcommand{\ncolor}[1]{\textcolor{blue}{#1}}  % Nathakhun's colored text
\newcommand{\mcolor}[1]{\textcolor{red}{#1}}  % Wasamon's colored text
\newcommand{\ncomm}[1]{\marginpar{\tiny\ncolor{#1}}}

%%

\newcommand{\rcolor}[1]{\textcolor{red}{#1}}  
\newcommand{\rcomm}[1]{\marginpar{\tiny\rcolor{#1}}}


\newcommand{\todoA}[1]{\todo[color=yellow!50]{#1}}


\newcommand{\bea}{\begin{eqnarray}}
\newcommand{\ena}{\end{eqnarray}}
\newcommand{\beas}{\begin{eqnarray*}}
\newcommand{\enas}{\end{eqnarray*}}
\newcommand{\beq}{\begin{equation}}
\newcommand{\enq}{\end{equation}}
\def\qed{\hfill \mbox{\rule{0.5em}{0.5em}}}
\newcommand{\bbox}{\hfill $\Box$}
\newcommand{\ignore}[1]{}
\newcommand{\ignorex}[1]{#1}
%\newcommand{\ws}{{\widetilde \sigma}}
%\newcommand{\abs}[1]{\left\vert#1\right\vert}
\newcommand{\wtilde}[1]{\widetilde{#1}}
\newcommand{\qmq}[1]{\quad\mbox{#1}\quad}
\newcommand{\qm}[1]{\quad\mbox{#1}}
\newcommand{\nn}{\nonumber}
\newcommand{\Bvert}{\left\vert\vphantom{\frac{1}{1}}\right.}
\newcommand{\To}{\rightarrow}
\newcommand{\E}{\mathbb{E}}
\newcommand{\Var}{\mathrm{Var}}
\newcommand{\Cov}{\mathrm{Cov}}
\newcommand{\Corr}{\mathrm{Corr}}
\newcommand{\dist}{\mathrm{dist}}
\newcommand{\diam}{\mathrm{diam}}
\newtheorem{theorem}{Theorem}[section]
\newtheorem{corollary}[theorem]{Corollary}
\newtheorem{conjecture}[theorem]{Conjecture}
\newtheorem{proposition}[theorem]{Proposition}
\newtheorem{lemma}[theorem]{Lemma}
\newtheorem{definition}[theorem]{Definition}
\newtheorem{example}[theorem]{Example}
\newtheorem{remark}[theorem]{Remark}
\newtheorem{case}{Case}[section]
\newtheorem{condition}{Condition}[section]

%%%%

%%%%%=============================================================================%%%%
%%%%  Remarks: This template is provided to aid authors with the preparation
%%%%  of original research articles intended for submission to journals published 
%%%%  by Springer Nature. The guidance has been prepared in partnership with 
%%%%  production teams to conform to Springer Nature technical requirements. 
%%%%  Editorial and presentation requirements differ among journal portfolios and 
%%%%  research disciplines. You may find sections in this template are irrelevant 
%%%%  to your work and are empowered to omit any such section if allowed by the 
%%%%  journal you intend to submit to. The submission guidelines and policies 
%%%%  of the journal take precedence. A detailed User Manual is available in the 
%%%%  template package for technical guidance.
%%%%%=============================================================================%%%%

\begin{comment}
    

%% as per the requirement new theorem styles can be included as shown below
\theoremstyle{thmstyleone}%
\newtheorem{theorem}{Theorem}%  meant for continuous numbers
%%\newtheorem{theorem}{Theorem}[section]% meant for sectionwise numbers
%% optional argument [theorem] produces theorem numbering sequence instead of independent numbers for Proposition
\newtheorem{proposition}[theorem]{Proposition}% 
%%\newtheorem{proposition}{Proposition}% to get separate numbers for theorem and proposition etc.

\theoremstyle{thmstyletwo}%
\newtheorem{example}{Example}%
\newtheorem{remark}{Remark}%

\theoremstyle{thmstylethree}%
\newtheorem{definition}{Definition}%
\end{comment}

\raggedbottom
%%\unnumbered% uncomment this for unnumbered level heads

\begin{document}

\title[Article Title]{An approximate zero bias transformation for random sums: Applications to sampling with outliers, auto insurance, and generative AI}

%%=============================================================%%
%% GivenName	-> \fnm{Joergen W.}
%% Particle	-> \spfx{van der} -> surname prefix
%% FamilyName	-> \sur{Ploeg}
%% Suffix	-> \sfx{IV}
%% \author*[1,2]{\fnm{Joergen W.} \spfx{van der} \sur{Ploeg} 
%%  \sfx{IV}}\email{iauthor@gmail.com}
%%=============================================================%%


    

\author[1,3]{\fnm{Wasamon} \sur{Jantai}}\email{wasamon.j@chula.ac.th}
\equalcont{These authors contributed equally to this work.}


\author*[2,3]{\fnm{Nathakhun} \sur{Wiroonsri}}\email{nathakhun.wir@kmutt.ac.th}
\equalcont{These authors contributed equally to this work.}

\affil[1]{\orgdiv{Department of Mathematics
		and Computer Science}, \orgname{Chulalongkorn University}, \orgaddress{\state{Bangkok}, \country{Thailand}}}
\affil*[2]{\orgdiv{Department of Mathematics}, \orgname{King Mongkut's University of Technology Thonburi}, \orgaddress{\city{Bangkok}, \country{Thailand}}}

\affil[3]{Statistics, Probability, and Data Science with R programming (SPD$\epsilon$R) research group}


%%==================================%%
%% Sample for unstructured abstract %%
%%==================================%%

\abstract{
We develop $L^1$ bounds for the difference between a test function of a random sum and a standard normal random variable, where the summands are assumed to be independent but not necessarily identically distributed. The bounds are obtained through a new version of the approximate zero bias transformation specifically developed for random sums. Although the identical distribution assumption is relaxed, the bounds are of order $1/\sqrt{n}$, matching the order of existing bounds in the literature under the same distributional assumption on the number of summands. The main results are then applied to three real-world settings: random sums obtained from simple random sampling with outliers, total insurance claims, and generative AI response times.}

%%================================%%
%% Sample for structured abstract %%
%%================================%%

% \abstract{\textbf{Purpose:} The abstract serves both as a general introduction to the topic and as a brief, non-technical summary of the main results and their implications. The abstract must not include subheadings (unless expressly permitted in the journal's Instructions to Authors), equations or citations. As a guide the abstract should not exceed 200 words. Most journals do not set a hard limit however authors are advised to check the author instructions for the journal they are submitting to.
% 
% \textbf{Methods:} The abstract serves both as a general introduction to the topic and as a brief, non-technical summary of the main results and their implications. The abstract must not include subheadings (unless expressly permitted in the journal's Instructions to Authors), equations or citations. As a guide the abstract should not exceed 200 words. Most journals do not set a hard limit however authors are advised to check the author instructions for the journal they are submitting to.
% 
% \textbf{Results:} The abstract serves both as a general introduction to the topic and as a brief, non-technical summary of the main results and their implications. The abstract must not include subheadings (unless expressly permitted in the journal's Instructions to Authors), equations or citations. As a guide the abstract should not exceed 200 words. Most journals do not set a hard limit however authors are advised to check the author instructions for the journal they are submitting to.
% 
% \textbf{Conclusion:} The abstract serves both as a general introduction to the topic and as a brief, non-technical summary of the main results and their implications. The abstract must not include subheadings (unless expressly permitted in the journal's Instructions to Authors), equations or citations. As a guide the abstract should not exceed 200 words. Most journals do not set a hard limit however authors are advised to check the author instructions for the journal they are submitting to.}

\keywords{Artificial Intelligence, Central Limit Theorem, Normal approximation, Stein's method, zero biasing }

%%\pacs[JEL Classification]{D8, H51}

\pacs[MSC Classification]{60F05}

\maketitle

\section{Introduction}\label{sec:introduction}

A random sum is a sum in which the number of summands is itself a random variable. It is written in the form
\bea \label{rs}
W = \sum_{i=1}^N X_i,
\ena
where $X_i$ are any random variables, and $N$ is a non-negative integer-valued random variable. In this work, we assume $X_i$ to be independent but are not necessarily identical, and $N$ is independent of all $X_i$. The random sum is very useful in modeling real-world quantities, such as total claim amount, aggregate losses, transaction volumes, total infections over a fixed time period, and total wait time of users or customers. The number of summands is random for these situations \citep{gnedenko2020random, YonghintJantai2025}.  

Distributional approximation of a random sum has been studied for decades \citep{robbins1948asymptotic,renyi1963central,sunklodas2015normal}.
%\cite{robbins1948asymptotic,renyi1963central,sunklodas2015normal}. 
Most of them assume that $X_i$ are identical. This assumption is a limitation for real-world modeling. For instance, in auto insurance, if the number of claims reaches a certain upper bound, the company may be stricter in paying the full amount requested by the auto service provider, leading to a new random claim amount. In a digital platform, wait time will be longer if the number of concurrent users exceeds a certain amount. In this work, we plan to relax the identical assumption and develop normal approximation methods for a random sum of independent nonidentical random variables using Stein's method based on the approximate zero biasing technique \citep{NW2017}. This technique has an extended application in \cite{JW25}.  

The main results are then applied to three real-world examples: sampling from a population with outliers, insurance claims, and generative AI response time. The most well-regarded sampling technique is simple random sampling. It is known that the standardized average always converges to the standard normal distribution. However, it is sometimes known that the population of interest contains outliers, which causes the number of non-outliers to be random \citep{barnett1994outliers}. The number of insurance claims is known to be random and normally modeled by a Poisson distribution %(see e.g. Klugman et al. \citeyear{klugman2012loss}; Denuit et al. \citeyear{denuit2007actuarial}). 
\citep{klugman2012loss,denuit2007actuarial}.
The claim amounts are typically assumed to be independent and identically distributed. However, the distribution of the claim amount may change during a year due to seasonality and the current number of claims. Generative AIs such as ChatGPT, Gemini, and Claude are now among the most popular tools for everyone. Their uses range from asking silly questions to helping code very difficult tasks. The wait time for answers depends on how simple the question is and how many concurrent users there are during the period. Though this exact same question has not been studied before in a probability and statistics context, statistical models used in Generative AI are summarized in \cite{dobriban2025statistical}.  

Stein's method was first introduced by Charles Stein in his seminal paper \citep{Stein72} and has become one of the most powerful methods to prove convergence in distribution. Its main advantages are that it provides non-asymptotic bounds on its target functions, which include distance and any other types, and that it can handle several situations involving dependence. For more details about the method in general, see the text \citep{CGS11} and the introductory notes \citep{Ross11}. The use of Stein's method for the normal approximation of random sums has been studied for more than 20 years \citep{barbour2006normal,chaidee2008berry}. More recently, \cite{Dal22} obtained an $L^1$ bound of order $O(1/\sqrt{n})$ under the assumption that the $X_i$ are independent and identically distributed.







\begin{comment}
\mcolor{The concept of Stein's method, particularly for normal approximation, is based on the fact that a random variable $W$ has the standard normal distribution, denoted $\mathcal{N}(0,1)$, if and only if
	\beas
	\E W f(W) = \E f'(W)
	\enas
	for all absolutely continuous functions $f$ with $\E |f'(W)| < \infty$. This equation, along with the form of the distance in \eqref{Wasdef}, leads to the differential equation
	\bea \label{steineq}
	h(w) - N h = f'_h(w)-wf_h(w),
	\ena
	where $Nh = \E h(Z)$ with $Z \sim \mathcal{N}(0,1)$, $h \in \mathcal{H}_1$, and $f_h$ is the solution to \eqref{steineq}. Taking the supremum over all $h \in \mathcal{H}_1$   in the expectation on the left-hand side of \eqref{steineq}, with $w$ replaced by a variable $W$, yields the distance between $W$ and $Z$ in \eqref{Wasdef}. Thus, one can handle the expectation on the right-hand side using the bounded solution $f_h$ of \eqref{steineq}, thereby avoiding direct computations of the distances.}

\end{comment}
%%%%%%%%%%


We recall that the $L^1$  distance between the distributions ${\cal L}(X)$ and ${\cal L}(Y)$ of real valued random variables $X$ and $Y$ is given, respectively, by
\bea \label{Wasdef}
d_1 \big({\cal L}(X),{\cal L}(Y)\big) &=& \int_{-\infty}^\infty |P(X \le t) - P(Y \le t)| dt \nn \\
&=& \sup_{h \in \mathcal{H}_1} |\E h(X)-\E h(Y)| 
\ena 
where $\mathcal{H}_1 =\{h: |h(y)-h(x)| \le |y-x|\}$. As previously mentioned, we provide an $L^1$ bound between a random sum given in \eqref{rs} and the standard normal, where we allow $X_i$ to be nonidentical by using an approximate zero bias transformation introduced in  \cite{NW2017}. We note that, as part of our generalization, the original zero biasing technique  \citep{GR97} used in \cite{EKJ09} is not applicable in our setting. 

We first recall the definition of a zero bias distribution. For a mean zero random variable $W$ with $\Var(W) = \sigma^2 < \infty$, $W^z$ has the zero-bias distribution with respect to $W$ if 
\bea \label{zerodef}
\E Wf(W) = \sigma^2\E f'(W^z)
\ena
for all absolutely continuous functions $f$ such that $\E |Wf(W)|<\infty$. The distribution of $W^z$ is known to be unique, and $W^z=W$ when $W$ is normally distributed. \cite{GR97} also presented a construction of a zero bias coupling for a sum of independent mean zero random variables by randomly selecting a summand based on its variance and replacing the selected summand by its zero bias variable. The approximate zero bias variable for the sum of independent mean zero random variables that we introduce in the next section is a variable $W^*$ defined similarly to $W^z$ with a remainder term added to the right hand side of \eqref{zerodef}.   



The remainder of this work is organized as follows. The approximate zero biasing structure for random sums is introduced in Section \ref{main} along with the approximation results. Then the applications and simulations mentioned previously are presented in Section \ref{sec:app}. Section \ref{sec:sum} is devoted to discussion and conclusion.



\section{Approximate Zero Bias Transformation and Main Bounds} \label{main}


In this section, we introduce an approximate zero bias transformation for a random sum of independent mean zero random variables with random upper limit and develop a bound of the different between a test function of the sum and a standard normal based on this transformation.

\cite{GR97} presented a construction of a zero bias coupling for a sum of independent mean zero random variables by randomly selecting a summand based on its variance and replacing the selected summand by its zero bias variable.   

In this work, we generalize the idea of this construction with the sum replaced by a random sum. Let $X_1,X_2,\ldots$ be independent mean zero random variables such that $\Var(X_i) = \sigma_i^2$ and $N$ be a discrete random variable with support being a subset of $\mathbb{N}_0$ and independent of all $X_i$. Here, with an integer $k \ge 0$, we denote $\mathbb{N}_k$ as the set of positive integers starting at $k$. Now consider $W$ as given in \eqref{rs}.
It is not hard to show that $\sigma^2 := \Var(W) = \sum_{i=1}^\infty\sigma_i^2 P(N\ge i)$. We propose the following construction.

\textbf{
	\begin{flushleft}
		Construction of an approximate zero bias variable:
	\end{flushleft}
} \label{construction}
\begin{enumerate}
	\item For $i\in [n]$, let $X_i^*$ has the zero bias distribution of $X_i$ independent of $(X_j)_{j\ne i}$ and $(X_j^*)_{j\ne i}$.
	\item Choose a random summand $X_I$ with 
	\bea \label{Idef}
	P(I=i) = \frac{\sigma_i^2P(N'\ge i)}{\sigma^2},
	\ena  
	where $N'$ has the same distribution as $N$ and independent of $N$ and $I$ is independent of all other random variables.
	\item Define 
	\bea \label{Wsdef}
	W^* = \begin{cases}
		\sum_{j\ne I}X_j+X_I^*, &\text{ \ if \ } N \ge I,\\
		W,  &\text{ \ if \ } N<I.
	\end{cases}
	\ena
\end{enumerate}

Next we show that the construction above leads to a variable which is similar to a zero bias variable with an additional remainder term. We say the variable $W^*$ in \eqref{approxzero} has an \textsl{approximate zero bias distribution} with respect to $W$. It is clear that an approximate zero bias variable is not unique, as we can always make an adjustment to the remainder term and get a new approximate zero bias variable. However, the key idea of this approximate version is that we need to construct a variable in such a way that the remainder term is sufficiently small. The concept of the approximate zero bias coupling was first introduced in \cite{NW2017}, where the variable was constructed via an approximate Stein's coupling.     

\begin{proposition} \label{2.1}
	Let $X_1,X_2,\ldots$ be independent mean zero random variables such that $\Var(X_i) = \sigma_i^2$ and $W=\sum_{i=1}^NX_i$ with $N$ be a discrete random variable with support being subset of $\mathbb{N}_0$ and independent of all $X_i$. Then $W^*$, constructed as in \eqref{Wsdef}, satisfies
	\bea \label{approxzero}
	\E Wf(W) = \sigma^2\E f'(W^*) + \sigma^2\left(\E f'(W^*)-\E[f'(W)|N<I]\right)\frac{P(N<I)}{P(N\ge I)}.
	\ena
\end{proposition}



\proof
Letting $f$ satisfy all the conditions stated in \eqref{zerodef} and starting on the left hand side of the equation, we have
\beas
\E Wf(W) &=& \E \left[\sum_{i=1}^N X_i f(W-X_i+X_i)\right] \\
&=& \sum_{k=0}^\infty \sum_{i=1}^k\E\left[ X_i f\left(\sum_{j=1}^kX_j-X_i+X_i\right)\right] P(N=k) \\
&=& \sum_{k=0}^\infty \sum_{i=1}^k \sigma_i^2 \E\left[ f'\left(\sum_{j=1}^kX_j-X_i+X_i^*\right)\right] P(N=k)  \\
&=& \sum_{i=1}^\infty \sigma_i^2 \sum_{k=i}^\infty \E\left[  f'\left(\sum_{j=1}^kX_j-X_i+X_i^*\right)\right] P(N=k) \\
&=& \sum_{i=1}^\infty \sigma_i^2 P(N\ge i) \sum_{k=i}^\infty \E\left[  f'\left(\sum_{j=1}^kX_j-X_i+X_i^*\right)\right] P(N=k|N\ge i) \\
&=& \sum_{i=1}^\infty \sigma^2  \E\left[  f'\left(\sum_{j=1}^NX_j-X_i+X_i^*\right)\Bigg|N\ge i\right]P(I=i)  \\
&=& \sigma^2 \E\left[f'(W-X_I+X_I^*)|N\ge I\right] = \sigma^2 \E\left[f'(W^*)|N\ge I\right].
\enas

Next, we handle the right hand side of \eqref{zerodef}. We have

\beas
\sigma^2\E[f'(W^*)] = \sigma^2\E\left[f'(W^*)|N\ge I\right]P(N\ge I)+\sigma^2\E\left[f'(W)|N< I\right]P(N< I)
\enas

It follows that
\beas
\E Wf(W) &=& \sigma^2\E f'(W^*) + \sigma^2\E f'(W^*)\left(\frac{1-P(N\ge I)}{P(N\ge I)}\right)-\sigma^2\E[f'(W)|N<I]\frac{P(N<I)}{P(N\ge I)} \\
&=&  \sigma^2\E f'(W^*) + \sigma^2\left(\E f'(W^*)-\E[f'(W)|N<I]\right)\frac{P(N<I)}{P(N\ge I)} .
\enas

\bbox



\begin{remark}
	Let $g(x) = f(x/\sigma)$. Then $g'(x) = f'(x/\sigma)/\sigma$. If we standardize $W$ to $W/\sigma$ to have variance one first, we have
	\bea \label{approxzero2}
	&& \E\left[\frac{W}{\sigma}f\left(\frac{W}{\sigma}\right)\right] = \frac{1}{\sigma}\E[Wg(W)] \nn \\ 
	&=& \sigma \E g'(W^*) + \sigma(\E g'(W^*)-\E[g'(W)|N<I])\frac{P(N<I)}{P(N\ge I)} \nn\\
	&=&  \E f'\left(\frac{W^*}{\sigma}\right) + \left(\E f'\left(\frac{W^*}{\sigma}\right) -\E \left[f'\left(\frac{W^*}{\sigma}\right) \Big|N<I \right]\right) \frac{P(N<I)}{P(N\ge I)} .
	\ena
\end{remark}


Now we provide a bound for the remainder term in \eqref{approxzero}. 

\begin{proposition}\label{prop:upperbound}
	Let $N$ be a discrete random variable with support being a subset of $\mathbb{N}_0$ and $\E N = \lambda$, and $I$ be as in \eqref{Idef}. Assume that $0<\sigma_i^2<\infty$ for all $i$. Then
	\bea \label{remainbound1}
	P(N < I) \le \frac{2\pi^2 \max_i \sigma_i^2\sqrt{\Var(N)}}{3\lambda\min_i\sigma_i^2}.
	\ena
	In particular, if $N$ is Poisson with mean $\lambda$, then
	\bea \label{remainbound2}
	P(N < I) \le \frac{2\pi^2\max_i \sigma_i^2 \sqrt{\lambda}}{3\lambda\min_i\sigma_i^2}.
	\ena
\end{proposition}
\proof
Since
\beas
P(N\ge I) &=& \sum_{n=1}^\infty \sum_{i=1}^n P(N=n)P(I=i) \\
&=& \sum_{n=1}^\infty \sum_{i=1}^n P(N=n)\frac{\sigma_i^2}{\sigma^2}P(N' \ge i) \\
&=& \sum_{i=1}^\infty \frac{\sigma_i^2}{\sigma^2}P(N' \ge i) \sum_{n=i}^\infty P(N=n) \\
&=& \sum_{i=1}^\infty \frac{\sigma_i^2}{\sigma^2}P^2(N \ge i), 
\enas
we have
\beas
P(N<I) &=& 1-\sum_{i=1}^\infty \frac{\sigma_i^2}{\sigma^2}P^2(N \ge i)  \\
&=& \frac{1}{\sigma^2}\sum_{i=1}^\infty \sigma_i^2 (P(N\ge i) - P^2(N \ge i))  \\
&\le& \frac{\max_i\sigma_i^2}{\sum_{j=1}^\infty \sigma_j^2 P(N\ge j)}\sum_{i=1}^\infty P(N \ge i)P(N<i)\\
&=& \frac{1}{\lambda}\frac{\max_i \sigma_i^2}{\min_i\sigma_i^2 } \sum_{i=1}^\infty P(N \ge i)P(N<i).
\enas
Using Chebyshev's inequality, we know that for $\E N - (k+1)\sqrt{\Var(N)} < i \le \E N - k\sqrt{\Var(N)}$,
\beas
P(N < i ) \le \frac{1}{k^2} \text{ \ and thus \ } P(N \ge i)P(N<i) \le \frac{1}{k^2},
\enas
and for $\E N + k\sqrt{\Var(N)} \le i < \E N + (k+1)\sqrt{\Var(N)}$,
\beas
P(N \ge i ) \le \frac{1}{k^2} \text{ \ and thus \ } P(N \ge i)P(N<i) \le \frac{1}{k^2}.
\enas
Therefore,
\beas
P(N<I) \le \frac{1}{\lambda}\frac{\max_i \sigma_i^2}{\min_i\sigma_i^2 }  \sum_{k=1}^\infty \frac{4\sqrt{\Var(N)}}{k^2}  = \frac{2\pi^2 \max_i \sigma_i^2\sqrt{\Var(N)}}{3\lambda\min_i\sigma_i^2}.
\enas

\bbox
\begin{comment}
	The following remark discusses a sufficient condition that the bound presented above converges to zero as the mean goes to infinity. 
	
	\begin{remark}
		The bound in \eqref{remainbound1} for $P(N< I)$ converges to zero as $\lambda \rightarrow \infty$ if $\Var(N) = o(\lambda^2)$. In the case that $N$ is Poisson, we have the bound in \eqref{remainbound2} for $P(N< I)$ be of order $\lambda^{-1/2}$.
	\end{remark}
	
	Next, we follow the idea in \cite{Elk09} to provide some interesting estimations. 
	
	\begin{proposition} \label{2.4} \mcolor{Let $X_1,X_2,\ldots$ be independent mean zero random variables such that $\Var(X_i) = \sigma_i^2$ and $W=\sum_{i=1}^NX_i$ with $N$ be a discrete random variable with support being subset of $\mathbb{N}_0$. Let $I$ be as in \eqref{Idef}.  Then} 
		\beas
		\E[X_I^*] = \frac{1}{2\sigma^2}\sum_{i=1}^\infty\E[X_i^3]P(N\ge i), \ \ \ \E\left[(X_I^*)^2\right] = \frac{1}{3\sigma^2}\sum_{i=1}^\infty\E[X_i^4]P(N\ge i) 
		\enas
		and
		\beas
		\E|W^*-W|^k = \frac{1}{2(k+1)\sigma^2} \sum_{i=1}^\infty\E\left[|X_i^s|^{k+2}\right]P^2(N\ge i),
		\enas
		where $X_i^s = X_i-\tilde{X_i}$ and $\tilde{X_i}$ is an independent copy of $X_i$.
	\end{proposition}
	
	\proof
	
	Using the zero bias condition with $g(x) = x^2$, we have
	\beas
	\E[X_I^*] &=& \sum_{i=1}^\infty\E[X_i^*]\frac{\sigma_i^2P(N\ge i)}{\sigma^2} \\
	&=& \sum_{i=1}^\infty\E[X_i^3]\frac{P(N\ge i)}{2\sigma^2} \\
	&=& \frac{1}{2\sigma^2}\sum_{i=1}^\infty\E[X_i^3]P(N\ge i).
	\enas
	
	The proof of the second equality is similar to the first one with $g(x)=x^3$. Next, we prove the last equality in the statement of the theorem.
	
	\beas
	\E|W^*-W|^k &=& \E\left[|W^*-W|^k \mathbf{1}(N\ge I) \right]+\E\left[|W^*-W|^k\mathbf{1}(N< I)\right]\\
	&=& \E\left[|X_I^*-X_I|^k \mathbf{1}(N\ge I) \right] \\
	&=& \sum_{i=1}^\infty\E\left[|X_i^*-X_i|^k \mathbf{1}(N\ge i) \right] \frac{\sigma_i^2P(N'\ge i)}{\sigma^2} \\
	&=& \sum_{i=1}^\infty\E\left[|X_i^*-X_i|^k \right]P(N\ge i) \frac{\sigma_i^2P(N'\ge i)}{\sigma^2} \\
	&=& \frac{1}{2(k+1)\sigma^2}\sum_{i=1}^\infty\E\left[|X_i^s|^{k+2} \right]P^2(N\ge i) .
	\enas
	
	\bbox
\end{comment}
%%%%%%%%%%%%%%%%%%%%%%%%%%%%%%%%%%%%%%%%%%%%%%%%%%%%%%%%%%%%%%%%%%%%%%%%%%%%%%%%%

The following theorem and its corollary provide bounds for \eqref{Wasdef}.

\begin{theorem}\label{uniform}
	Let $X_1,X_2,\ldots$ be independent mean zero random variables such that $\Var(X_i) = \sigma_i^2$. Let $W=\sum_{i=1}^NX_i$, where $N$ is a discrete random variable with support being a subset of $\mathbb{N}_0$ and independent of all $X_i$.
		Define $\sigma^2 := \Var(W) = \sum_{i=1}^\infty\sigma_i^2 P(N\ge i)$. Let $I$ be as in \eqref{Idef}.  Then
	\begin{align*}
		\left|\E h\left(\frac{W}{\sigma}\right)-N(h)\right|\leq \frac{||f'_h||}{\sigma^3}\sum_{i=1}^\infty \left({1\over 2}\E|X_i|^3+ \sigma_i^2\E|X_i| \right)P^2(N\ge i)+2 ||f^{'}_{h}|| \cdot {P(N < I)},
	\end{align*}
	where $N(h) = \E h(Z)$ with $Z \sim \mathcal{N}(0,1)$, $h \in \mathcal{H}_1$ and $f_h$ is the solution to the Stein's equation.
	\end{theorem}

    
	\proof 
	By Stein's equation and Proposition 2.1, we have
	\begin{align*}
&\left|\E h\left(\frac{W}{\sigma}\right)-N(h)\right| = \left|\E\left(f_{h}'(\frac{W}{\sigma})-(\frac{W}{\sigma})f_{h}(\frac{W}{\sigma})\right)\right|\\
&\leq \left|\E\left(f_{h}'(\frac{W}{\sigma})-f_{h}'(\frac{W^{*}}{\sigma})\right)\right|+\left| \E\left(f_{h}'\left(\frac{W^{*}}{\sigma}\right) \right)-\E\left(f_{h}'\left(\frac{W}{\sigma}\right)\mid N<I \right)\right|\cdot\displaystyle\frac{P(N<I)}{P(N\geq I)}\\
& \leq  \left|\E\left(f_{h}'(\frac{W}{\sigma})-f_{h}'(\frac{W^{*}}{\sigma})\right)\right|+\left| \E\left(f_{h}'\left(\frac{W^{*}}{\sigma}\right) \right)-\E\left(f_{h}'(\frac{W}{\sigma})\mid N<I \right)\right|\cdot\displaystyle\frac{P(N<I)}{P(N\geq I)}\\
&\leq  \frac{1}{\sigma} ||f_{h}^{''}||\E \left|W^{*}-W\right|+\left| \E\left(f_{h}'(\frac{W^{*}}{\sigma}) \right)-\E\left(f_{h}'(\frac{W}{\sigma})\mid N<I \right)\right|\cdot\displaystyle\frac{P(N<I)}{P(N\geq I)}.
\end{align*}
%By 
%the property that $||f_{k}^{''}||\leq 2||h_{k}^{'}||$ and 
%Proposition \ref{2.4} we have,
%\begin{align*}
%\left|\E\left(f_{h}'(\frac{W}{\sigma})-f_{h}'(\frac{W^{*}}{\sigma}\right)\right|\leq &\frac{1}{\sigma} ||f_{h}^{''}||\E \left|W^{*}-W\right|
%% \leq   \frac{1}{4\sigma^{3}}||f_{h}^{''}||\sum^{\infty}_{i=1}\E\left|X_{i}^{s}\right|^{3}P^{2}(N\geq i).
%\end{align*}
Notice that 
\beas
\E|W^*-W|&=& \E\left[|W^*-W| \mathbf{1}(N\ge I) \right]+\E\left[|W^*-W|\mathbf{1}(N< I)\right]\\
&=& \E\left[|X_I^*-X_I| \mathbf{1}(N\ge I) \right] \\
&=& \sum_{i=1}^\infty\E\left[|X_i^*-X_i| \mathbf{1}(N\ge i) \right] \frac{\sigma_i^2P(N'\ge i)}{\sigma^2} \\
&=& \sum_{i=1}^\infty\E|X_i^*-X_i| \,P(N\ge i) \frac{\sigma_i^2P(N'\ge i)}{\sigma^2} \\
&=& \frac{1}{\sigma^2}\sum_{i=1}^\infty \sigma_i^2\E|X_i^*-X_i| P^2(N\ge i)\\
&\le & \frac{1}{\sigma^2}\sum_{i=1}^\infty \sigma_i^2\left(\E|X_i^*|+\E|X_i| \right)P^2(N\ge i)\\
&\le & \frac{1}{\sigma^2}\sum_{i=1}^\infty \left({1\over 2 }\E|X_i|^3+ \sigma_i^2\E|X_i| \right)P^2(N\ge i)
\enas
and  
\beas
&&\left| \E\left(f_{h}'\left(\frac{W^{*}}{\sigma}\right) \right)-\E\left[f_{h}'\left(\frac{W}{\sigma}\right)\mid N<I \right]\right|\cdot\displaystyle\frac{P(N<I)}{P(N\geq I)}\\
&=&\Big| \E\left[f_{h}'\left(\frac{W^{*}}{\sigma}\right) \mid N<I\right]P(N<I)+\E\left[f_{h}'\left(\frac{W^{*}}{\sigma}\right) \mid N\ge I\right]P(N\ge I)\\
&&\quad\quad\quad\quad\quad\quad\quad\quad\quad\quad\quad\quad\quad\quad-\E\left[f_{h}'(\frac{W}{\sigma})\mid N<I \right]\Big|\cdot\displaystyle\frac{P(N<I)}{P(N\geq I)}\\
&=&\Big| \E\left[f_{h}'\left(\frac{W^{*}}{\sigma}\right) -f_{h}'\left(\frac{W}{\sigma}\right) \mid N<I\right]-\E\left[f_{h}'(\frac{W^*}{\sigma})\mid N<I \right]P(N\ge I)\\
&&\quad\quad\quad\quad\quad\quad+\E\left[f_{h}'(\frac{W^{*}}{\sigma}) | N\ge I\right]P(N\ge I)\Big|\cdot\displaystyle\frac{P(N<I)}{P(N\geq I)}\\
&=&\Big| \E\left[f_{h}'(\frac{W^{*}}{\sigma}) | N\ge I\right]P(N\ge I)-\E\left[f_{h}'(\frac{W^*}{\sigma})\mid N<I \right]P(N\ge I)
\Big|\cdot\displaystyle\frac{P(N<I)}{P(N\geq I)}\\
&\le & 2||f'_h||P(N<I).
\enas
Therefore, 
\begin{align*}
\left|\E h\left(\frac{W}{\sigma}\right)-N(h)\right|\leq  \frac{||f'_h||}{\sigma^3}\sum_{i=1}^\infty \left({1\over 2}\E|X_i|^3+ \sigma_i^2\E|X_i| \right)P^2(N\ge i)+2 ||f^{'}_{h}|| \cdot{P(N < I)}.
\end{align*}
\bbox

\begin{corollary} \label{coro2.5} Let $X_1,X_2,\ldots$ be independent mean zero random variables such that $\Var(X_i) = \sigma_i^2$. Let $W=\sum_{i=1}^NX_i$, where $N$ is  a discrete random variable with support being subset of $\mathbb{N}_0$, independent of $\{X_i\}_{i=1}^\infty$.
	Define $\sigma^2 := \Var(W) = \sum_{i=1}^\infty\sigma_i^2 P(N\ge i)$. Let $I$ be as in \eqref{Idef}. Then 
	$$d_1\left(\mathcal{L}(W/\sigma), \mathcal{L}(Z) \right)\le \frac{1}{\sigma^3}\sum_{i=1}^\infty \left(\E|X_i|^3+ 2 \sigma_i^2\E|X_i| \right)P^2(N\ge i)+2 \sqrt{\frac{2}{\pi}}\cdot\displaystyle P(N<I),$$  where $Z$ is a standard normal random variable.
\end{corollary}

The following remark gives a special case for i.i.d. random variables and a binomial random variable $N$.
\begin{remark} \label{Remark:BN} Consider i.i.d. mean zero random variables $X_1,X_2, \dots$  such that $\Var(X_i) = \sigma^2_1$ for all $i$. Define $W$ as in \eqref{rs}, where $N$ is  a discrete random variable with support being subset of $\mathbb{N}_0$ and independent of all $X_i$, independent of $\{X_i\}_{i=1}^\infty$. Let $I$ be as in \eqref{Idef} and \\${\sigma}^2 := \Var(W) = \sum_{i=1}^\infty\sigma_i^2 P(N\ge i)=\sigma^2_1\E[N]$. If $N\sim Bin(n,p)$ has a binomial distribution with parameters $n$ and probability $p$, then, by Corollary \ref{coro2.5}, 

\begin{equation}  \label{Theoretical UB}
    d_1\left(\mathcal{L}(W/{\sigma}), \mathcal{L}(Z) \right)= \left({\E|X_1|^3 \over (\sigma^2_1 np)^{3/2}} +{2\sigma_1^2\E|X_1|\over (\sigma^2_1np)^{3/2}}\right)\sum_{i=1}^n P^2(N\ge i)  +2 \sqrt{\frac{2}{\pi}}\cdot\displaystyle P(N<I),\end{equation} 
    where $Z$ is a standard normal random variable.
\begin{comment}
\begin{equation}  \label{Theoretical UB}
    d_1\left(\mathcal{L}(W/{\sigma}), \mathcal{L}(Z) \right)= \left({\E|X_1|^3 \over (\sigma^2_1 np)^{3/2}} +{2\sigma_1^2\E|X_1|\over (\sigma^2_1np)^{3/2}}\right)\sum_{i=1}^n P^2(N\ge i)  +{2\sqrt{2} \over \sqrt{\pi}}\cdot\displaystyle{\left(1-{1\over np}\sum_{i=1}^nP^2(N\ge i)\right)},\end{equation} \todo{Modify}
    where $Z$ is a standard normal random variable.
\end{comment}

In comparison, \cite{Dal22} showed that 
\begin{equation} \label{Daly's Theoretical UB}
d_1\left(\mathcal{L}(W/{\sigma}), \mathcal{L}(Z) \right)\le \frac{2}{\sqrt{np\sigma^2_1}}\left(\frac{\E|X_1|^3}{2\sigma^2_1}+p\E|X_1| \right).
\end{equation}

It follows immediately from \eqref{Theoretical UB} that the first term is $O(1/\sqrt{n})$. By Proposition \ref{prop:upperbound}, the second term is also $O(1/\sqrt{n})$. Hence, the two terms have the same convergence rate. Although the bounds in \eqref{Theoretical UB} and \eqref{Daly's Theoretical UB} are both of order $O(1/\sqrt{n})$, the constant of  \eqref{Daly's Theoretical UB} is actually smaller. We show a numerical example in Table \ref{tab:daly_comparison} to demonstrate it. However, our bound in \eqref{Theoretical UB} does not require the identical assumption, and it will benefit users as illustrated through the applications in Section \ref{sec:app}.  
\end{remark}

\newpage
\begin{table}[htbp]
\centering
\begin{tabular}{|c|c|c|c|c|c|}
\hline
\multirow{2}{*}{$n$}
& Upper bounds
& \multirow{2}{*}{Daly's upper bounds}
& \multirow{2}{*}{$n$}
& Upper bounds
& \multirow{2}{*}{Daly's upper bounds}
\\
& from \eqref{Theoretical UB}
&
&
& from \eqref{Theoretical UB}
&
\\
\hline
$2^1$  & 3.2140 & 3.9315 & $2^{11}$ & 0.1919 & 0.1229 \\
$2^2$  & 2.8350 & 2.7800 &$2^{12}$ & 0.1364 & 0.0869 \\
$2^3$  & 2.3179 & 1.9658 &$2^{13}$ & 0.0968 & 0.0614 \\
$2^4$  & 1.8025 & 1.3900 &$2^{14}$ & 0.0686 & 0.0434 \\
$2^5$  & 1.3579 & 0.9829 &$2^{15}$ & 0.0486 & 0.0307 \\
$2^6$  & 1.0021 & 0.6950 &$2^{16}$ & 0.0344 & 0.0217 \\
$2^7$  & 0.7296 & 0.4914 &$2^{17}$ & 0.0243 & 0.0154 \\
$2^8$  & 0.5264 & 0.3475 &$2^{18}$ & 0.0172 & 0.0109 \\
$2^9$  & 0.3775 & 0.2457 &$2^{19}$ & 0.0122 & 0.0077 \\
$2^{10}$ & 0.2696 & 0.1738 &$2^{20}$ & 0.0086 & 0.0054 \\
\hline
\end{tabular}
\caption{Comparison of the upper bounds derived from  \eqref{Theoretical UB} and  \eqref{Daly's Theoretical UB}. The upper bounds are computed from \( X_i = Y_i - \mathbb{E}[Y_i] \), \( Y_i \sim \mathrm{Gamma}(k = 0.8203, \theta = 10962.0161) \), 
and \( N \sim {Bin}(n, 2583/8753) \).}
\label{tab:daly_comparison}
\end{table}





\section{Applications} \label{sec:app}

In this section, we demonstrate real-life applications of our theoretical results, including random sampling with outliers, insurance claims, and generative AI response times.


\subsection{Simple random sampling with outliers}

\begin{comment}
The non-outlier elements are i.i.d.\ copies of a centered exponential random variable, denoted by
\[
X_1, \dots, X_N \overset{\text{i.i.d.}}{\sim} A,
\]
while the outliers are i.i.d.\ copies of a different distribution,
\[
X_{N+1}, \dots, X_{N+m} \overset{\text{i.i.d.}}{\sim} O.
\]
The details of \(A\) and \(O\) are provided in the next subsection.
Moreover, the non-outlier sequence and the outlier sequence are assumed to be independent.    
\end{comment}

In this subsection, we apply our main results to simple random sampling with outliers. The traditional zero bias transformation does not work in this case due to extra terms caused by outliers.

\subsubsection{Theoretical results}
Let $m,N \in \mathbb{N}$ such that $m \ll N$. We consider a finite population of size $N+m$,
\[
\mathcal{A}=\{x_1,\dots,x_N,\,x_{N+1},\dots,x_{N+m}\},
\]
consisting of $N$ non-outliers and $m$ outliers. A simple random sample of size $n \in \mathbb{N}$ is drawn without replacement from $\mathcal{A}$. Let $M$, independent of all $X_i$, be the number of non-outliers in the sample. Then $M$ has a hypergeometric distribution, denoted by,
\[
M \sim H(N+m,\, n,\, N), \text{ and }
\mathbb{E}[M] = \frac{nN}{N+m}.
\]
The support of $M$ is $\max\{0,n-m\} \le M\le \min\{n,N\}$. After reordering the sample, we write
\[
\mathbf{X} = (X_1,\dots,X_M, X_{M+1},\dots,X_n),
\]
where the first $M$ components are non-outliers and the remaining $n-M$ are outliers.

\begin{comment}
 \[
X_1,\dots,X_M \stackrel{d}{=} A,
\qquad
X_{M+1},\dots,X_p \stackrel{d}{=} O.
\]   
\end{comment}

Let
\[
W = \sum_{i=1}^{M} X_i.
\]
Then, denoting $\Var(X_i) = s^2$, the variance of $W$ is given by 	
\begin{align}\label{SRS:sigma}
\sigma^2=s^2\sum_{i=1}^\infty P(M\ge i)=s^2\E M={n Ns^2\over N+m}=O(n).
\end{align}



The following corollary provides an upper bound for the Wasserstein distance between $W/\sigma$ and a standard normal random variable. 


\begin{corollary} \label{SRM:Theoretical result} Let $X_1,X_2,\ldots$ be i.i.d. mean zero random variables such that $\Var(X_i) = s^2$ for all $i\geq 1$. Let $W=\sum_{i=1}^MX_i$, where $M$ is  a hypergeometric random variable, independent of $\{X_i\}_{i=1}^\infty$, and let $\sigma^2$ be defined as in  \eqref{SRS:sigma}. Then 
	\begin{align*} d_1\left(\mathcal{L}(W/\sigma), \mathcal{L}(Z) \right)  \le &\left({\E|X_1|^3 \over \left({nN s^2\over N+m}\right)^{3/2}} +{2s^2\E|X_1|\over \left({nNs ^2\over N+m}\right)^{3/2}}\right)\sum_{i=1}^{\min\{n,N\}} P^2(M\ge i)  \\&\quad \quad+{2\sqrt{2} \over \sqrt{\pi}}\cdot\displaystyle{\left(1- \left( \frac{N+m}{nN}\right)\sum_{i=1}^{\min\{n,N\}} P^2(M\ge i)\right)}, \end{align*}  where $Z$ is a standard normal random variable.
\end{corollary}
\proof Following the argument in the proof of Proposition \ref{prop:upperbound} and using \eqref{SRS:sigma}, we have
\begin{align*}
		P(M \ge I) &=  \sum_{i=1}^{\min\{n,N\}} \frac{\Var(X_i)}{\sigma^2} \, P^2(M \ge i)\\
		&=  \sum_{i=1}^{\min\{n,N\}} \frac{s^2}{\sigma^2} \, P^2(M \ge i)\\& =  \left( \frac{N+m}{nN}\right)\sum_{i=1}^{\min\{n,N\}} \, P^2(M \ge i).
	\end{align*} 
	Using Theorem~\ref{uniform} together with the bounds 
	$\|f_h''\| \le 2$ and $\|f_h'\| \le \sqrt{2/\pi}$, 
	we obtain \begin{align*} d_1\left(\mathcal{L}(W/\sigma), \mathcal{L}(Z) \right) & \le \left({\E|X_1|^3 \over \left({nNs^2\over N+m}\right)^{3/2}} +{2s^2\E|X_1|\over \left({nNs^2\over N+m}\right)^{3/2}}\right)\sum_{i=1}^{\min\{n,N\}} P^2(M\ge i)  \\&\quad \quad+{2\sqrt{2} \over \sqrt{\pi}}\cdot\displaystyle{\left(1- \left( \frac{N+m}{nN}\right)\sum_{i=1}^{\min\{n,N\}} P^2(M\ge i)\right)}.\end{align*}

\bbox




\subsubsection{Simulation} 
We now apply Corollary~\ref{SRM:Theoretical result} to an exponential population. More precisely, let $Y_1,Y_2,\ldots$ be i.i.d.\ exponential random
variables with mean $\lambda$, and define the centered variables
\[
X_i = Y_i-\lambda.
\]
For $X_i\sim \mathrm{Exp}(1/\lambda)$, note that $\lambda^{-1} X_i\sim \mathrm{Exp}(1)$. 
Note that the corresponding first three absolute central moments are as follows:
\[
	\E|X_i-\E X_i| = \frac{2\lambda}{ e}, \qquad
    \E|X_i-\E X_i|^2 = {\lambda^2}, \qquad
    \E|X_i-\E X_i|^3 = {\lambda^3}\left(\frac{12}{e}-2\right).
	\]
    Substituting these quantities into Corollary~\ref{SRM:Theoretical result}, and noting that $n \ll N$  in our setting, yields the following bound:
    \begin{align} d_1\left(\mathcal{L}(W/\sigma), \mathcal{L}(Z) \right) \le &\left({\left( \dfrac{12}{e}-2\right) \over \left({nN\over N+m}\right)^{3/2}} +{\left( \dfrac{4}{e}\right)\over \left({nN\over N+m}\right)^{3/2}}\right)\sum_{i=1}^{n} P^2(M\ge i) \nonumber  \\&\quad \quad+{2\sqrt{2} \over \sqrt{\pi}}\cdot\displaystyle{\left(1- \left( \frac{N+m}{nN}\right)\sum_{i=1}^{n} P^2(M\ge i)\right)} \label{SRM:simulation result},\end{align} where $Z$ is a standard normal random variable. The bound is independent of the exponential mean parameter $\lambda$. 
    

\begin{comment}
    

\ncolor{***The corollary must be in general for any population. (not specific to the exponential population.) Then, in the simulation section, we consider an Exponential population and we will plug in the first three moments there.***}
\end{comment}
 
\begin{comment}
\begin{corollary}\label{Cor-outliers1}
	Let $W$ be a mean zero, variance $\sigma^2 >0$ random variable, then
	\begin{align*} &d_1\left(\mathcal{L}(W/\sigma), \mathcal{L}(Z) \right) \\&\quad \le\left({\left( \dfrac{12}{e}-2\right) \over \left({nN\over N+m}\right)^{3/2}} +{\left( \dfrac{4}{e}\right)\over \left({nN\over N+m}\right)^{3/2}}\right)\sum_{i=\max\{0,n-m\}}^{\min\{n,N\}} P^2(M\ge i)  \\&\quad \quad+{2\sqrt{2} \over \sqrt{\pi}}\cdot\displaystyle{\left(1- \left( \frac{N+m}{nN}\right)\sum_{i=\max\{0,n-m\}}^{\min\{n,N\}} P^2(M\ge i)\right)},\end{align*} where $Z$ is a standard normal random variable.
\end{corollary}
\begin{proof}
	Using Theorem~\ref{uniform} together with the bounds 
	$\|f_h''\| \le 2$ and $\|f_h'\| \le \sqrt{2/\pi}$, and the identity
	\begin{align*}
		P(M \ge I) 
		&= \sum_{i=\max\{0,n-m\}}^{\min\{n,N\}}  \frac{\sigma_i^2}{\sigma^2} \, P(M \ge i)^2\\&=  \sum_{i=\max\{0,n-m\}}^{\min\{n,N\}} \frac{r^2}{\sigma^2} \, P(M \ge i)^2\\& =  \left( \frac{N+m}{nN}\right)\sum_{i=\max\{0,n-m\}}^{\min\{n,N\}} \, P(M \ge i)^2,
	\end{align*} 
	we obtain \begin{align*} &d_1\left(\mathcal{L}(W/\sigma), \mathcal{L}(Z) \right) \\&\quad \le \left({\E|X_1|^3 \over \left({nNr^2\over N+m}\right)^{3/2}} +{2r^2\E|X_1|\over \left({nNr^2\over N+m}\right)^{3/2}}\right)\sum_{i=\max\{0,n-m\}}^{\min\{n,N\}} P^2(M\ge i)  \\&\quad \quad+{2\sqrt{2} \over \sqrt{\pi}}\cdot\displaystyle{\left(1- \left( \frac{N+m}{nN}\right)\sum_{i=\max\{0,n-m\}}^{\min\{n,N\}} P^2(M\ge i)\right)}.\end{align*}
	For $X\sim \mathrm{Exp}(\lambda)$, note that $\lambda X\sim \mathrm{Exp}(1)$. Using the identities
	\[
	\E|X-\E X|^3 = \frac{1}{\lambda^3}\left(\frac{12}{e}-2\right), \qquad
	\E|X-\E X|^2 = \frac{1}{\lambda^2}, \qquad
	\E|X-\E X| = \frac{2}{\lambda e},
	\]
	the desired expression follows.
\end{proof}
\end{comment}
Now we perform a simulation to see how our bound in \eqref{SRM:simulation result} behaves compared to the real distance between the normalized $W$ and $Z$. We consider a population with the non-outlier size of \(N = 100{,}000{,}000\), and the outlier size of \(m = 1{,}200{,}000\). The nonoutliers are constructed from an Exponential distribution with mean $5$. The outliers are created from mixed Uniform distributions,
\[
O \sim 
\begin{cases}
	\mathrm{Unif}(-20,-10), & \text{with probability } p_0=\frac{55}{70},\\[0.4em]
	\mathrm{Unif}(50,60), & \text{with probability } 1-p_0.
\end{cases}
\]
Note that $\mathbb{E}[O]=0$.

\begin{comment}
Since the population is large and $X_1,\ldots,X_n$ are a simple random sample, it is safe to say approximately that 
\beas
X_i \sim Exp(1/5) \text{ for } i=1,\ldots,M,
\enas
and 
\beas
X_i \sim O \text{ for } i=M+1,\ldots,n.
\enas

\end{comment}

Table \ref{tab:samp} shows the upper bounds from \eqref{SRM:simulation result} and the experimental distance computed from 100,000 simulated samples.  

\begin{table}[!htbp]
		\centering
	\begin{tabular}{|c|c|c|c|c|c|}
		\hline
		$n$ & Upper bound  & Simulated distance
		& $n$ & Upper bound & Simulated distance   \\
		\hline
		$2^1$   & 2.75061143 & 0.19881889 & $2^{9}$  & 0.17661563 & 0.01431534 \\
		$2^2$   & 1.95054948 & 0.15390773 & $2^{10}$ & 0.12499765 & 0.01214080 \\
		$2^3$   & 1.38449178 & 0.11342520 & $2^{11}$ & 0.08844060 & 0.00853899 \\
		$2^4$   & 0.98356093 & 0.08105740 & $2^{12}$ & 0.06256320 & 0.00800496 \\
		$2^5$   & 0.69896369 & 0.05700365 & $2^{13}$ & 0.04425172 & 0.00759138 \\
		$2^6$   & 0.49638711 & 0.04091913 & $2^{14}$ & 0.03129700 & 0.00484885 \\
		$2^7$   & 0.35206696 & 0.02958365 & $2^{15}$ & 0.02213342 & 0.00400173 \\
		$2^8$   & 0.24944274 & 0.01988771 & $2^{16}$ & 0.01565219 & 0.00877013 \\
		\hline
	\end{tabular}
	\caption{The comparison of theoretical upper bounds and simulated distances for different values of $n$. The simulated distance is computed from 100,000 repeated simulations.}
    \label{tab:samp}
\end{table}
\begin{comment}
\begin{table}[!htbp]
		\centering
	\begin{tabular}{|c|c|c|c|c|c|}
		\hline
		$n$ & Upper bound  & Simulated distance
		& $n$ & Upper bound & Simulated distance   \\
		\hline
		$2^1$   & 3.34189070 & 0.19881889 & $2^{9}$  & 0.17380298 & 0.01431534 \\
		$2^2$   & 2.04134862 & 0.15390773 & $2^{10}$ & 0.12354132 & 0.01214080 \\
		$2^3$   & 1.35746929 & 0.11342520 & $2^{11}$ & 0.08769475 & 0.00853899 \\
		$2^4$   & 0.94444462 & 0.08105740 & $2^{12}$ & 0.06218402 & 0.00800496 \\
		$2^5$   & 0.67035278 & 0.05700365 & $2^{13}$ & 0.04405992 & 0.00759138 \\
		$2^6$   & 0.47888103 & 0.04091913 & $2^{14}$ & 0.03120032 & 0.00484885 \\
		$2^7$   & 0.34218232 & 0.02958365 & $2^{15}$ & 0.02208480 & 0.00400173 \\
		$2^8$   & 0.24410034 & 0.01988771 & $2^{16}$ & 0.01562778 & 0.00877013 \\
		\hline
	\end{tabular}
	\caption{The comparison of theoretical upper bounds and simulated distances for different values of $n$. The simulated distance is computed from 100,000 repeated simulations.}
    \label{tab:samp}
\end{table}
\todoA{fix table}
\end{comment}
\begin{figure}[!htbp]
	\centering
    \includegraphics[width=0.8\textwidth]{Fig1eee.png}
	\caption{Simulated and theoretical upper bounds as functions of $n$, plotted on a logarithmic scale with a consistent $x$-axis spacing.}
	\label{fig:placeholder}
\end{figure}


Table \ref{tab:samp} and Figure \ref{fig:placeholder} illustrate the comparison between the theoretical upper bounds (blue) and the simulated distances between $W$ and the standard normal distribution for different values of $n$ ranging from $2$ to $2^{16}$. The results show that the simulated distances are always smaller than the theoretical upper bounds for all values of $n$. The convergence rates of the two are aligned, although our theoretical upper bound converges much more slowly due to the large constant.



\newpage
\subsection{Auto Insurance claims}

In this subsection, we approximate the number of auto insurance claims on our generated synthetic dataset by a normal distribution. We generate claim amounts using a Gamma distribution and the total number of claims using a Binomial distribution. The distributional parameters are estimated by the maximum likelihood (MLE) from publicly available data on Kaggle \citep{Xiaomeng}. 
\begin{comment}\todo{Have to mention data in detail. Sample size. Number of claims. Number of policies.}\end{comment} 


\begin{comment}
The dataset used in this study is synthetically generated to model the total amount of claims. The number of claims follows a Binomial distribution, while each individual claim is drawn from a Gamma distribution and adjusted by subtracting its mean to ensure a mean of zero. The total amount of the claim is then calculated as the sum of these adjusted claims. The simulation was implemented in R. 


We use publicly available data from Kaggle, specifically the ``Car Insurance Claim Data'' dataset by Xiaomeng Sun (https://www.kaggle.com/datasets/xiaomengsun/car-insurance-claim-data). This data set contains information on car insurance claims, including factors such as policyholder details, claim amounts, and other relevant insurance-related factors.
\end{comment}


\subsubsection{Theoretical results}
We consider i.i.d. mean-zero random variables $X_1,X_2,\ldots$ with finite variance. Let $N$ be a binomial random variable independent of all $X_i$ with parameters $n$ and $p$, the number of policyholders and the fraction of policyholders who make at least one claim, respectively. Let
\beas
W=\sum_{i=1}^NX_i.
\enas 

 
 Again, following the argument in the proof of Proposition \ref{prop:upperbound}, we have
 
 $$\sigma^2 := \Var(W) = \sum_{i=1}^\infty \Var(X_i) P(N\ge i)=\sum_{i=1}^\infty\E X^2_i P(N\ge i)=\Var(X_1)\E[N]=\Var(X_1) np,$$
and  \[
P(N\ge I) = \sum_{i=1}^\infty \frac{\Var(X_i)}{\sigma^2}P^2(N \ge i) ={1\over np}\sum_{i=1}^nP^2(N\ge i).
\]

By Remark \ref{Remark:BN}, the bound on the Wasserstein distance between $W$ and the standard normal can be rewritten as follows: \begin{align}
    \label{Remark-insurance}d_1\left(\mathcal{L}(W/{\sigma}), \mathcal{L}(Z) \right)\le \frac{1}{\sigma^3}\sum_{i=1}^{n} \left(\E|X_1|^3 + 2\E X_1^2 \E|X_1|\right) P^2(N \ge i)  +{2\sqrt{2} \over \sqrt{\pi}}\cdot\displaystyle{\left(1-{1\over np}\sum_{i=1}^nP^2(N\ge i)\right)},
    \end{align} where $Z$ is a standard normal random variable.

\begin{comment}
    
\rcolor{The bound of the Wasserstein distance between $W$ and the standard normal is presented in the corollary below, following from Corollary \ref{coro2.5}.} 
\todoA{This is the same as Remark 2.6 (see Equation (8)), so we should remove one of them.}
\begin{corollary}\label{Cor-insurance}  Let $X_1,X_2,\ldots$ be i.i.d. mean zero random variables with finite variance. Let $W=\sum_{i=1}^N X_i$, where $N$ is  a binomial random variable  with parameters $n$ and $p$, independent of $\{X_i\}_{i=1}^\infty$.
	Define $\sigma^2 := \Var(W) = \sum_{i=1}^\infty Var(X_i) P(N\ge i)$. Let $I$ be as in \eqref{Idef}. Then
			$$d_1\left(\mathcal{L}(W/{\sigma}), \mathcal{L}(Z) \right)= \frac{1}{\sigma^3}\sum_{i=1}^{n} \left(\E|X_1|^3 + 2\E X_1^2 \E|X_1|\right) P^2(N \ge i)  +{2\sqrt{2} \over \sqrt{\pi}}\cdot\displaystyle{\left(1-{1\over np}\sum_{i=1}^nP^2(N\ge i)\right)},$$ where $Z$ is a standard normal random variable.
		\end{corollary}
        \begin{proof}
The result follows immediately from
Corollary~\ref{coro2.5}.
\end{proof}
\end{comment}


\subsubsection{Simulation} 


The original dataset \citep{Xiaomeng} consists of 10,302 records corresponding to 8,753 unique policyholder IDs. Among these records, 2,746 have positive claim amounts. After aggregating records with the same policyholder ID, the final dataset contains 8,753 policyholders, of whom 2,583 have a positive total claim amount and 6,170 have a total claim amount of zero. The Gamma parameters are estimated by MLE using all positive claims, yielding $k=0.8203$ and $\theta=10962.0161$. The fraction of policyholders who make at least one claim is $\frac{2583}{8753}$, which gives the MLE of $p$.

Specifically, we consider independent gamma random variables $X'_1,X'_2,\ldots$ with shape parameter $k=0.8203$ %$k=0.820285415321782 $
and scale parameter $\theta=10962.0161$. Let $N$ be a binomial random variable independent of all $X'_i$ with parameters $n$ and $p=\frac{2583}{8753}$, where $p$ represents the fraction of policyholders who make at least one claim Let
\beas
W=\sum_{i=1}^NX_i,
\enas 
where $X_i=X'_i-\E[X'_i]$ for all $i$. Then $$\sigma^2 := \Var(W) = \sum_{i=1}^\infty Var(X_i) P(N\ge i)=\sum_{i=1}^\infty Var( X'^2_i )P(N\ge i)=k\theta^2\E[N]=k\theta^2 np,$$
and  \[
P(N\ge I) = \sum_{i=1}^\infty \frac{Var(X_i)}{\sigma^2}P^2(N \ge i) = \sum_{i=1}^\infty \frac{Var(X'_i)}{\sigma^2}P^2(N \ge i)={1\over np}\sum_{i=1}^nP^2(N\ge i).
\]
For the specified distribution of $X_i$, we substitute these quantities into \eqref{Remark-insurance} to obtain the theoretical upper bounds shown in Table \ref{tab:insure}. We compute the first three absolute central moments numerically, since their analytical expressions are rather complicated. We also perform a simulation study to compare these theoretical bounds with the simulated distance between the normalized $W$ and $Z$. We repeat the simulation 1,000,000 times and compute the average distance from the standard normal distribution, as described in the theoretical part above.


\begin{comment}
Substituting these quantities into \eqref{Remark-insurance} yields the following  bound:
			\begin{align} \label{eqn:single population} d_1\left(\mathcal{L}(W/{\sigma}), \mathcal{L}(Z) \right)= \frac{1}{\sigma^3}\sum_{i=1}^{n} \left(\E|X_1|^3 + 2\E X_1^2 \E|X_1|\right) P^2(N \ge i)  +{2\sqrt{2} \over \sqrt{\pi}}\cdot\displaystyle{\left(1-{1\over np}\sum_{i=1}^nP^2(N\ge i)\right)}, \end{align} where $Z$ is a standard normal random variable.
	\todoA{This part needs to be revised.}
			
       Now we perform a simulation to see how our bound in \eqref{eqn:single population} behaves compared to the real distance between the normalized $W$ and $Z$. 
\end{comment}
\newpage
\begin{table}[htbp]
			\centering
			\begin{tabular}{|c|c|c|c|c|c|}
				\hline
				$n$ &  Upper bounds & Simulated distance & $n$ &  Upper bounds & Simulated distance\\
				\hline
				$2^{1}$  & 3.182252156 & 0.397469296  & $2^{11}$ & 0.189724807 &  0.014580690\\
				$2^{2}$  & 2.805325953 & 0.270339191  & $2^{12}$ & 0.134803390 & 0.010225892\\
				$2^{3}$  & 2.292840787 & 0.216659633 & $2^{13}$ & 0.095644132 & 0.007348849 \\
				$2^{4}$  & 1.782629262 & 0.160307504 & $2^{14}$ & 0.067792444 & 0.005204517 \\
				$2^{5}$  & 1.342652002 & 0.115341608 & $2^{15}$ & 0.048017396 & 0.003791629 \\
				$2^{6}$  & 0.990774915 & 0.081930908 & $2^{16}$ & 0.033993871 & 0.002766599 \\
				$2^{7}$  & 0.721327391 & 0.058025642 & $2^{17}$ & 0.024057517 & 0.002042665  \\
				$2^{8}$  & 0.520433949 & 0.040943521 & $2^{18}$ & 0.017021343 & 0.001708912 \\
				$2^{9}$  & 0.373190694 & 0.029063853 & $2^{19}$ & 0.012040961 & 0.001505429 \\
				$2^{10}$ & 0.266478699 & 0.020438395 & $2^{20}$ & 0.008516772 & 0.001407022 \\
				\hline
			\end{tabular}
            
			\caption{The comparison of theoretical upper bounds and simulated distances for different values of $n$. The simulated distance is computed from 1,000,000 repeated simulations.}
        \label{tab:insure}
		\end{table}
Figure \ref{fig:insure} compares the simulated distances with the theoretical upper bounds. The theoretical upper bounds (blue) decrease as $n$ increases, while the simulated values (red) are consistently smaller than the theoretical bounds. The gap between the two curves also becomes smaller, indicating that the approximation improves as $n$ increases. This improvement is due to the diminishing effect of the large constant in the theoretical bound.
		\begin{figure}[!htbp]  % 'h' means here, but you can use 't' for top, 'b' for bottom, etc.
			\centering
            \includegraphics[width=0.8\textwidth]{Fig2.png}
            \caption{Simulated and theoretical upper bounds as functions of $n$, plotted on a logarithmic scale with a consistent $x$-axis spacing.}
            \label{fig:insure}
		\end{figure}
		
		
		%%%%%%%%%%%%%%%
\newpage
		\subsection{Extension of auto insurance claims} \label{app}

In this subsection, we extend the analysis from the previous section to approximate insurance claims by a normal distribution under the additional assumption that ``an insurance company becomes stricter in approving claim amounts after the total claim amount reaches a certain level.'' This assumption is motivated by the hypothesis that auto service providers may submit slightly inflated bills to insurance companies. When the total claim amount reaches a certain level, the insurer may review the claim more carefully and negotiate a lower payment.
		
		\subsubsection{Theoretical results}

Let $n_1\in\mathbb{N}$. Consider two independent groups of mean zero random variables,
$X_1,\ldots,X_{n_1}$
and
$X_{n_1+1},X_{n_1+2},\ldots$,
each having finite variance. The random variables within each group are i.i.d.,
while the distributions of the two groups are allowed to differ.  Let $N$ be a binomial random variable with parameters
$n$ and $p$, independent of all the $X_i$, where $n$ is the number of policyholders, and $p$ represents the
fraction of policyholders who make at least one claim. Define 
		\[
		W = \sum_{i=1}^N X_i.
		\]

        
Similar to the previous case, we have
\begin{align}\label{TH:IC1}
			\sigma^2 := \Var(W) &= \sum_{i=1}^\infty \Var(X_i) \, P(N \ge i)\nn \\
			&=  \sum_{i=1}^{n_{1}  } \Var(X_1)  \, P(N \ge i)+ \sum_{i={n_{1}+1}  }^n \Var(X_{n_{1} + 1})  \, P(N \ge i),
		\end{align}
		and
		\begin{align}\label{TH:IC2}
			{P}(N \ge I) &= \sum_{i=1}^\infty \frac{\Var(X_i)}{\sigma^2}P^2(N \ge i) \nn\\
			&= \sum_{i=1}^{n_1}\frac{\Var(X_1)}{\sigma^2} \, {P}^2(N \ge i) +\sum_{i=n_1+1}^{n}\frac{\Var(X_{n_1+1})}{\sigma^2} \, {P}^2(N \ge i).
		\end{align}
		
		%\[
		%{P}(N \ge I) = \sum_{i=1}^\infty \frac{\sigma_i^2}{\sigma^2} \, {P}^2(N \ge i) = \frac{1}{np} \sum_{i=1}^n {P}^2(N \ge i).
		%\]
	\begin{remark}
If $n_1\ge n$, the setting reduces to the single population case.
	\end{remark}	
		
		The following corollary provides an upper bound for the Wasserstein distance between 
		$W/\sigma$ and a standard normal random variable. 
		
		\begin{corollary}\label{cor:insurance} Let $n_1\in \mathbb{N}.$ Let $X_1,X_2,\ldots$ be mean zero random variables with finite variances.
Suppose that
$X_1,\ldots,X_{n_1} \text{ and }
X_{n_1+1}, X_{n_1+1},\ldots$
form two independent groups of i.i.d.\ random variables, where the distributions may differ
between the two groups. Let $W=\sum_{i=1}^N X_i$, where $N$ is  a binomial random variable  with parameters $n$ and $p$, independent of $\{X_i\}_{i=1}^\infty$, and let $\sigma^2$ be defined as in \eqref{TH:IC1}.  %Let $I$ be as in \eqref{Idef}.
   Then
			\begin{align*}
				d_1\left(\mathcal{L}(W/\sigma), \mathcal{L}(Z) \right)
				&\le \frac{1}{\sigma^3}\sum_{i=1}^{n_{1}} \left(\E|X_1|^3 + 2\E X_1^2 \E|X_1|\right) P^2(N \ge i) \\
				&\quad + \frac{1}{\sigma^3} \sum_{i=n_{1}+1}^{n} \left(\E|X_{n_{1}+1}|^3 + 2\E X_{n_{1}+1}^2  \E|X_{n_{1}+1}|\right) P^2(N \ge i) \\
				&\quad + \frac{2\sqrt{2}}{\sqrt{\pi}} \left(1 - \sum_{i=1}^{n_1}\frac{\Var(X_1)}{\sigma^2} \, {P}^2(N \ge i) -\sum_{i=n_1+1}^{n}\frac{\Var(X_{n_1+1})}{\sigma^2} \, {P}^2(N \ge i)\right),
			\end{align*}
			where $Z$ is a standard normal random variable.
		\end{corollary}
		\proof The proof follows from Theorem~\ref{uniform}, together with the bounds
			$\|f_h''\| \le 2$ and $\|f_h'\| \le \sqrt{2/\pi}$, and equations \eqref{TH:IC1} and \eqref{TH:IC2}.
           \bbox
		
		
		\subsubsection{Simulation}   
        
		First we let \( N \sim \mathrm{Bin}(n, p) \), where \( p = \dfrac{2583}{8753} \) as in the previous section. Consider a sequence of independent claim amounts \( X'_1, X'_2, \ldots \), where the first $n_1$ variables, \( X'_1, \ldots, X'_{n_1} \), follow a Gamma distribution with shape parameter \( k_1 = 4.04 \) and scale parameter \( \theta_1 = 1025.78 \), and the remaining variables \( X'_{n_1+1}, X'_{n_1+2}, \ldots \) follow a Gamma distribution with parameters \( k_2 = 2.5 \) and \( \theta_2 = 800 \). These parameters are chosen for illustration so that the mean claim amount in the first group is larger than that in the second group based on our hypothesis. Each claim is centered by subtracting its expected value, that is, \( X_i = X'_i - \mathbb{E}[X'_i] \). Finally, define the total (centered) claim amount as
		\[
		W = \sum_{i=1}^N X_i.
		\]
		Then, the variance of \( W \) is
		\begin{align*}
			\sigma^2 := \Var(W) 
			&=  \sum_{i=1}^{n_{1}  } k_1\theta^2_1  \, P(N \ge i)+ \sum_{i={n_{1}+1}  }^n k_2\theta^2_2  \, P(N \ge i) .
		\end{align*}
		
		Moreover,
		
		\begin{align*}
			{P}(N \ge I) 
			&= \sum_{i=1}^{n_1}\frac{k_1\theta^2_1}{\sigma^2} \, {P}^2(N \ge i) +\sum_{i=n_1+1}^{n}\frac{k_2\theta^2_2}{\sigma^2} \, {P}^2(N \ge i).
		\end{align*}
Substituting these quantities into Corollary~\ref{cor:insurance} yields the following  bound:
			\begin{align}
				d_1\left(\mathcal{L}(W/\sigma), \mathcal{L}(Z) \right)
				&\le \frac{1}{\sigma^3}\sum_{i=1}^{n_{1}} \left(\E|X_1|^3 + 2\E X_1^2 \E|X_1|\right) P^2(N \ge i) \nonumber\\
				&\quad + \frac{1}{\sigma^3} \sum_{i=n_{1}+1}^{n} \left(\E|X_{n_{1}+1}|^3 + 2\E X_{n_{1}+1}^2  \E|X_{n_{1}+1}|\right) P^2(N \ge i)\nonumber \\
				&\quad + \frac{2\sqrt{2}}{\sqrt{\pi}} \left(1 - \sum_{i=1}^{n_1}\frac{k_1\theta^2_1}{\sigma^2} \, {P}^2(N \ge i) -\sum_{i=n_1+1}^{n}\frac{k_2\theta^2_2}{\sigma^2} \, {P}^2(N \ge i)\right). \label{eqn:two population}
			\end{align}
            
We now simulate the behavior of our bound in \eqref{eqn:two population} and compare it with the simulated distance between the normalized $W$ and $Z$, as shown in Table \ref{tab:insure2}. In \eqref{eqn:two population}, we again compute the first three absolute central moments numerically, since their analytical expressions are rather complicated.

			\begin{table}[!htbp]
			\centering
			\begin{tabular}{|c|c|c|c|c|c|}
				\hline
				$n$ & Upper bounds & Simulated distance & $n$ & Upper bounds & Simulated distance \\
				\hline
				$2^1$  & 2.837730952 & 0.336570427 & $2^{11}$ & 0.173259958& 0.008203795 \\
				$2^2$  & 2.483412912 & 0.199716890 & $2^{12}$ & 0.122745093 & 0.005945545 \\
				$2^3$  & 2.020426998 & 0.110711251 & $2^{13}$ & 0.086570289 & 0.004286396  \\
				$2^4$  & 1.566328648 & 0.077239348 & $2^{14}$ & 0.061084707 & 0.003213821 \\
				$2^5$  & 1.177558296 & 0.053436330 & $2^{15}$ & 0.043149835 & 0.002353975 \\
				$2^6$  & 0.867886841 & 0.037369233 & $2^{16}$ & 0.030503179 & 0.001809013 \\
				$2^7$  & 0.631338871 & 0.026270865 & $2^{17}$ & 0.021570680 & 0.001647794 \\
				$2^8$  & 0.455242470 & 0.018602045 & $2^{18}$ & 0.015255875 & 0.001461229 \\
				$2^9$  & 0.317432568 & 0.014071492 & $2^{19}$ & 0.010789914 & 0.001391717 \\
				$2^{10}$ & 0.239888497 & 0.011096590 & $2^{20}$ & 0.007631111 & 0.001401313 \\
				\hline
			\end{tabular}
			\caption{The comparison of theoretical upper bounds and simulated distances for different values of $n$. The simulated distance is computed from 1,000,000 repeated simulations.}
    \label{tab:insure2}
		\end{table}
		
		
	
		\newpage
	For visualization, we plot the simulated distances compared with the theoretical upper bounds in Figure \ref{fig:insure2}. The results behave similarly to the previous case.
    
		
		\begin{figure}[H]  
			\centering
            \includegraphics[width=0.8\textwidth]{Fig3.png}
            \caption{Simulated and theoretical upper bounds as functions of $n$, plotted on a logarithmic scale with a consistent $x$-axis spacing.}
        \label{fig:insure2}
		\end{figure}

		
		
		
		%%%%%%%%%%%%%%%%%%%%%%%%%%

			

		%%%%%%%%%%%%%%%%%%%%%%%%%%%%%%%%%%%%%%%%%%%%%
		
		\subsection{Generative AI} 

       In this subsection, we model the response time of a generative AI system under the hypothesis that the response time for the same question increases as user traffic increases. We use ChatGPT to support this hypothesis and to help define four appropriate levels of delay.
		\subsubsection{Theoretical results}
        \begin{comment}
            
    
		 Assume that many users submit moderate-complexity queries at the same time. 
		Based on publicly available information about ChatGPT’s large-scale system 
		(assisted by ChatGPT (OpenAI, January~15,~2026)), we divide system behavior into four states, 
		as shown in Table~\ref{table:concurrency_states}.
		
		\begin{table}[htbp] 
			\centering
			\begin{tabular}{|c|l|c|}
				\hline
				Concurrency level & System state & Selected value (used in simulation) \\ \hline
				1k--50k & No delay & 25k \\ \hline
				2k--100k & Noticeable but acceptable delay & 50k \\ \hline
				5k--200k & Severe delay (queue explosion) & 100k \\ \hline
				$>200$k & Throttling / rejection & -- \\ \hline
			\end{tabular}
			\caption{System states by concurrency level for moderate-complexity queries}
			\label{table:concurrency_states}
		\end{table}
		
		To describe different load levels, we use concurrency levels of 25,000, 50,000, and 100,000 for each group, from fast to very slow service (see Table \ref{table:concurrency_states}). 
		In the simulations and in computing the upper bounds from Corollary~\ref{Cor: ChatGPT}, we set
		\[
		m_1 = 0.2~\text{s}, \quad 
		m_2 = 1~\text{s}, \quad 
		m_3 = 5~\text{s}, \quad 
		m_4 = 20~\text{s},
		\]
		where \(m_j\) is the mean waiting time in group \(j\). These correspond to the rates
		\[
		\lambda_1 = 5, \quad 
		\lambda_2 = 1, \quad 
		\lambda_3 = 0.2, \quad 
		\lambda_4 = 0.05.
		\] These values are chosen to illustrate system behavior under different load conditions, where \(m_j\) is the mean waiting time in group \(j\).
		
		


        
		Let \( N \sim \mathrm{Poisson}(\lambda) \) be independent of the sequence \(\{X'_i\}_{i\geq 1}\).  
		Consider a sequence of waiting times \( X'_1, X'_2, \ldots \), where the population is divided into \( 4 \) groups according to indices  
		\[
		0 = n_0< n_1 < n_2< n_3.
		\]
		For each group \( j \in \{1, 2, 3\} \), the waiting times  
		\[
		X'_{n_{j-1}+1}, X'_{n_{j-1}+2}, \ldots, X'_{n_j}
		\]
		are independent and follow the Exponential distribution with rate parameter \( \lambda_j \).   For the last group \( j = 4 \), the range extends to infinity, that is,
		\[
		X'_{n_3+1}, X'_{n_3+2}, \ldots,
		\] and these waiting times follow \( \mathrm{Exp}(\lambda_4) \).
		%The parameters are specified as:
		%\[
		%\lambda_1 = 5, \quad 
		%\lambda_2 = 1, \quad 
		%\lambda_3 = 0.2, \quad 
		%\lambda_4 = 0.05,
		%\]
		%where each rate is measured in units of per second, so the waiting times are in seconds.
		For each group $j=1,2,3$ and indices $i=n_{j-1}+1,\ldots,n_j$, define
		\[
		X_i = X'_i - \frac{1}{\lambda_j}.
		\]
		Finally, define the total centered waiting time as
		\[
		W = \sum_{i=1}^{N} X_i.
		\]
  \end{comment}
         Let $n_1, n_2, n_3\in \mathbb{N}$ be such that $n_1< n_2< n_3$. Let $X_1,X_2,\ldots$ be mean zero random variables with finite variances.
Suppose that
$X_1,\ldots,X_{n_1}$, $X_{n_1+1},\ldots X_{n_2}$, $X_{n_2+1},\ldots X_{n_3}$    and  $X_{n_3+1}, X_{n_3+2}, \ldots $ 
form four independent groups of i.i.d.\ random variables, where the distributions may differ among the four groups. Let $W=\sum_{i=1}^N X_i$, where $N$ is a Poisson random variable  with parameter $\lambda$, independent of $\{X_i\}_{i=1}^\infty$.


We again follow the proof of Proposition \ref{prop:upperbound} to obtain the following:
		\begin{align} \label{TH:AI1}
			\sigma^2 := \Var(W) &= \sum_{i=1}^\infty \Var(X_i) \, P(N \ge i) \nonumber\\
			&=  \sum_{i=1}^{n_{1}  } \Var(X_1)  \, P(N \ge i)+ \sum_{i={n_{1}+1}  }^{n_2}\Var(X_{n_{1} + 1})  \, P(N \ge i)\nonumber \\
			&\quad+  \sum_{i={n_2+1}}^{n_{3}  }\Var(X_{n_{2} + 1}) \, P(N \ge i)+  \sum_{i=n_{3}+1}^{\infty  }\Var(X_{n_{3} + 1})  \, P(N \ge i). 
		\end{align}
		Moreover,
		\begin{align} \label{TH:AI2}
				P(N \ge I) 
				&= \sum_{i=1}^{n_1}\frac{\Var(X_1)}{\sigma^2} \, {P}^2(N \ge i) +\sum_{i=n_1+1}^{n_2}\frac{\Var(X_{n_1+1})}{\sigma^2} \, {P}^2(N \ge i)\nonumber\\
                &+\sum_{i=n_2+1}^{n_3}\frac{\Var(X_{n_2+1})}{\sigma^2} \, {P}^2(N \ge i)+\sum_{i=n_3+1}^{\infty}\frac{\Var(X_{n_3+1})}{\sigma^2} \, {P}^2(N \ge i),
			\end{align} 
		
	The following corollary establishes an upper bound on the Wasserstein distance between $W/\sigma$ and the standard normal distribution.
		\begin{corollary} \label{Cor: ChatGPT} 
 Let $n_1, n_2, n_3\in \mathbb{N}$ be such that $n_1< n_2< n_3$. Let $X_1,X_2,\ldots$ be mean zero random variables with finite variances.
Suppose that
$X_1,\ldots,X_{n_1}$, $X_{n_1+1},\ldots X_{n_2}$, $X_{n_2+1},\ldots X_{n_3}$    and  $X_{n_3+1}, X_{n_3+2}, \ldots $ 
form four independent groups of i.i.d.\ random variables, where the distributions may differ among the four groups. Let $W=\sum_{i=1}^N X_i$, where $N$ is a Poisson random variable  with parameter $\lambda$, independent of $\{X_i\}_{i=1}^\infty$, and let
 $\sigma^2$ be defined as in \eqref{TH:AI1}. Then
			\begin{align}\label{Cor:Chatgt}
				d_1\!\left(\mathcal{L}(W/\sigma), \mathcal{L}(Z) \right)
				&\le \frac{1}{\sigma^3} \sum_{i=1}^{n_{1}} 
				\left(\mathbb{E}|X_1|^3 + 2\,\mathbb{E} X_1^2 \,\mathbb{E}|X_1|\right) P^2(N \ge i) \nn\\[4pt]
				&\quad + \frac{1}{\sigma^3} \sum_{i=n_{1}+1}^{n_2} 
				\left(\mathbb{E}|X_{n_{1}+1}|^3 + 2\,\mathbb{E} X_{n_{1}+1}^2  \,\mathbb{E}|X_{n_{1}+1}|\right) P^2(N \ge i)\nn \\[4pt]
				&\quad + \frac{1}{\sigma^3} \sum_{i=n_{2}+1}^{n_3} 
				\left(\mathbb{E}|X_{n_{2}+1}|^3 + 2\,\mathbb{E} X_{n_{2}+1}^2  \,\mathbb{E}|X_{n_{2}+1}|\right) P^2(N \ge i)\nn \\[4pt]
				&\quad + \frac{1}{\sigma^3} \sum_{i=n_{3}+1}^{\infty} 
				\left(\mathbb{E}|X_{n_{3}+1}|^3 + 2\,\mathbb{E} X_{n_{3}+1}^2  \,\mathbb{E}|X_{n_{3}+1}|\right) P^2(N \ge i) \nn\\[4pt]
				&\quad + \frac{2\sqrt{2}}{\sqrt{\pi}} 
				\Bigg(1 -\sum_{i=1}^{n_1}\frac{1}{\lambda^2_1\sigma^2} \, {P}^2(N \ge i) -\sum_{i=n_1+1}^{n_2}\frac{1}{\lambda^2_2\sigma^2} \, {P}^2(N \ge i)\nn\\&\quad\quad\quad\quad  -\sum_{i=n_2+1}^{n_3}\frac{1}{\lambda^2_3\sigma^2} \, {P}^2(N \ge i)-\sum_{i=n_3+1}^{\infty}\frac{1}{\lambda^2_4\sigma^2} \, {P}^2(N \ge i)\Bigg).
			\end{align}
		\end{corollary}
		\proof The proof follows from Theorem~\ref{uniform}, together with the bounds
			$\|f_h''\| \le 2$ and $\|f_h'\| \le \sqrt{2/\pi}$, and equations \eqref{TH:AI1} and \eqref{TH:AI2}. 
		\begin{comment}	Using  Theorem~\ref{uniform} together with the bounds 
			$\|f_h''\| \le 2$ and $\|f_h'\| \le \sqrt{2/\pi}$, and the identity
			\begin{align*}
				P(Z \ge I) 
				&= \sum_{i=1}^{n_1}\frac{1}{\lambda^2_1\sigma^2} \, {P}^2(N \ge i) +\sum_{i=n_1+1}^{n_2}\frac{1}{\lambda^2_2\sigma^2} \, {P}^2(N \ge i)\\&+\sum_{i=n_2+1}^{n_3}\frac{1}{\lambda^2_3\sigma^2} \, {P}^2(N \ge i)+\sum_{i=n_3+1}^{\infty}\frac{1}{\lambda^2_4\sigma^2} \, {P}^2(N \ge i),
			\end{align*} 
			we obtain 
			\begin{align*}
				d_1\!\left(\mathcal{L}(W/\sigma), \mathcal{L}(Z) \right)
				&\le \frac{1}{\sigma^3} \sum_{i=1}^{n_{1}} 
				\left(\mathbb{E}|X_1|^3 + 2\,\mathbb{E} X_1^2 \,\mathbb{E}|X_1|\right) P^2(N \ge i) \\[4pt]
				&\quad + \frac{1}{\sigma^3} \sum_{i=n_{1}+1}^{n_2} 
				\left(\mathbb{E}|X_{n_{1}+1}|^3 + 2\,\mathbb{E} X_{n_{1}+1}^2  \,\mathbb{E}|X_{n_{1}+1}|\right) P^2(N \ge i) \\[4pt]
				&\quad + \frac{1}{\sigma^3} \sum_{i=n_{2}+1}^{n_3} 
				\left(\mathbb{E}|X_{n_{2}+1}|^3 + 2\,\mathbb{E} X_{n_{2}+1}^2  \,\mathbb{E}|X_{n_{2}+1}|\right) P^2(N \ge i) \\[4pt]
				&\quad + \frac{1}{\sigma^3} \sum_{i=n_{3}+1}^{\infty} 
				\left(\mathbb{E}|X_{n_{3}+1}|^3 + 2\,\mathbb{E} X_{n_{3}+1}^2  \,\mathbb{E}|X_{n_{3}+1}|\right) P^2(N \ge i) \\[4pt]
				&\quad + \frac{2\sqrt{2}}{\sqrt{\pi}} 
				\Bigg(1 -\sum_{i=1}^{n_1}\frac{1}{\lambda^2_1\sigma^2} \, {P}^2(N \ge i) -\sum_{i=n_1+1}^{n_2}\frac{1}{\lambda^2_2\sigma^2} \, {P}^2(N \ge i)\\&\quad\quad\quad\quad  -\sum_{i=n_2+1}^{n_3}\frac{1}{\lambda^2_3\sigma^2} \, {P}^2(N \ge i)-\sum_{i=n_3+1}^{\infty}\frac{1}{\lambda^2_4\sigma^2} \, {P}^2(N \ge i)\Bigg).
			\end{align*}
			For $X\sim \mathrm{Exp}(\lambda)$, note that $\lambda X\sim \mathrm{Exp}(1)$. Using the identities
			\[
			\E|X-\E X|^3 = \frac{1}{\lambda^3}\left(\frac{12}{e}-2\right), \qquad
			\E|X-\E X|^2 = \frac{1}{\lambda^2}, \qquad
			\E|X-\E X| = \frac{2}{\lambda e},
			\]
			the desired expression follows. \bbox
            \end{comment}
	\bbox

        \begin{comment}
		\begin{remark}
			\rcolor{We remark that our bound in Corollary \ref{Cor: ChatGPT} has three jumps. When $\lambda$ is close to the threshold $n_1, n_2$ or $n_3$, $N$ is more likely to exceed the threshold, which results in more effects from the second to the fourth terms in \eqref{Cor:Chatgt}, and $1/\lambda_j$ is much larger than $1/\lambda_{j-1}$.}
		\end{remark}
	\todoA{Consider relocating this remark. It would also be helpful to state that $\lambda_j$ is decreasing in $j$.}	
		\end{comment}
		
		
		\subsubsection{Simulation} 
        
Assume that many users submit moderate-complexity queries at the same time. 
With the assistance of ChatGPT (OpenAI, January~15,~2026), we classify 
the system behavior into four states: no delay, noticeable but acceptable delay, 
severe delay (queue explosion), and throttling/rejection. We use these four states to model the system response time in both the theoretical analysis and the simulation.

Let \( N \sim \mathrm{Poisson}(\lambda) \) be independent of the sequence \(\{X'_i\}_{i\geq 1}\).  
		Consider a sequence of waiting times \( X'_1, X'_2, \ldots \),
        where the waiting times are divided into four groups according to indices  
		\[
		0 = n_0< n_1 < n_2< n_3.
		\]
		For each group \( j \in \{1, 2, 3\} \), the waiting times  
		\[
		X'_{n_{j-1}+1}, X'_{n_{j-1}+2}, \ldots, X'_{n_j}
		\]
		are independent and follow the Exponential distribution with rate parameter \( r_j \).   For the last group \( j = 4 \), the range extends to infinity.
	
		For each group $j=1,2,3,4$, define
		\[
		X_i = X'_i - \frac{1}{r_j}.
		\]
		Finally, define the total centered waiting time as
		\[
		W = \sum_{i=1}^{N} X_i.
		\]
For $X\sim \mathrm{Exp}(r)$, note that $r X\sim \mathrm{Exp}(1)$.  We then apply the preceding result by
substituting the corresponding first three absolute central moments:
			\[
			  \E|X-\E X| = \frac{2}{r e},
            \qquad
			\E|X-\E X|^2 = \frac{1}{r^2}, \qquad
            \E|X-\E X|^3 = \frac{1}{r^3}\left(\frac{12}{e}-2\right).
			\]
Substituting these quantities into Corollary~\ref{Cor: ChatGPT} yields the following  bound: \begin{align}\label{sim:AI1}
				d_1\!\left(\mathcal{L}(W/\sigma), \mathcal{L}(Z) \right)
				&\le \frac{1}{\sigma^3} \sum_{i=1}^{n_{1}} 
				\left(\frac{1}{r_1^3}\left(\frac{12}{e}-2\right) + \frac{4}{er_1^3 }\right) P^2(N \ge i) \nn \\[4pt] 
				&\quad + \frac{1}{\sigma^3} \sum_{i=n_{1}+1}^{n_2} 
				\left(\frac{1}{r_2^3}\left(\frac{12}{e}-2\right) + \frac{4}{er_2^3 }\right) P^2(N \ge i) \nn\\[4pt]
				&\quad + \frac{1}{\sigma^3} \sum_{i=n_{2}+1}^{n_3} 
				\left(\frac{1}{r_3^3}\left(\frac{12}{e}-2\right) + \frac{4}{er_3^3 }  \right) P^2(N \ge i)\nn \\[4pt]
				&\quad + \frac{1}{\sigma^3} \sum_{i=n_{3}+1}^{\infty} 
				\left(\frac{1}{r_4^3}\left(\frac{12}{e}-2\right) + \frac{4}{er_4^3 }  \right) P^2(N \ge i)\nn \\[4pt]
				&\quad + \frac{2\sqrt{2}}{\sqrt{\pi}} 
				\Bigg(1 -\sum_{i=1}^{n_1}\frac{1}{r^2_1\sigma^2} \, {P}^2(N \ge i) -\sum_{i=n_1+1}^{n_2}\frac{1}{r^2_2\sigma^2} \, {P}^2(N \ge i) \nn\\
				&\quad\quad\quad\quad-\sum_{i=n_2+1}^{n_3}\frac{1}{r^2_3\sigma^2} \, {P}^2(N \ge i)-\sum_{i=n_3+1}^{\infty}\frac{1}{r^2_4\sigma^2} \, {P}^2(N \ge i)\Bigg),
			\end{align}
			where $Z$ is a standard normal random variable.

        
		For the simulations and the computation of the upper bounds in \eqref{sim:AI1}, we set
		\begin{comment}\[
		m_1 = 0.2~\text{s}, \quad 
		m_2 = 1~\text{s}, \quad 
		m_3 = 5~\text{s}, \quad 
		m_4 = 20~\text{s},
		\]
		where \(m_j\) is the mean wait time in group \(j\) in second. These correspond to the rates \end{comment}
		\[
		r_1 = 5, \quad 
		r_2 = 1, \quad 
		r_3 = 0.2, \quad 
		r_4 = 0.05.
		\]
We choose $r_1>r_2>r_3>r_4$  so that the mean waiting time increases across the groups. We select threshold values 
$n_1 = 25{,}000$, $n_2 = 50{,}000$, and $n_3 = 100{,}000$ according to ChatGPT.
		\begin{comment}    
		Let \( N \sim \mathrm{Poisson}(\lambda) \), where
		\begin{multline*}
			\lambda \in \{2^k : 1 \le k \le 14\}
			\cup \{19384, 22384, 25384,  32768, 40168, 47568\} \\
			\cup \{54968, 65536, 85536\}
			\cup \{105536, 125536, 131072, 262144, 524288\}.
		\end{multline*}
		\todoA{Maybe adjust the way the values of $\lambda$ are presented.}
        \end{comment}
Let \(N \sim \mathrm{Poisson}(\lambda)\), where
\[
\lambda \in \{2^k : 1 \le k \le 19\}
\cup
\{19384, 22384, 25384, 40168, 47568, 54968,
85536, 105536, 125536\}.
\]
Now we perform a simulation to examine how our bound in \eqref{sim:AI1} behaves compared with the simulated distance between the normalized $W$ and $Z$ as shown in Table \ref{tab:AI}. Note that we intentionally include additional values of $\lambda$ such that the means of the Poisson distribution are around the thresholds, allowing us to examine how the bounds behave near these thresholds.
		
		
		\begin{table}[htbp]
			\centering
			\begin{tabular}{|c|c|c|c|c|c|}
				\hline
				$\lambda$ & Upper bounds & Simulated distance &  $\lambda$ & Upper bounds & Simulated distance \\
				\hline
				2      & 2.303442184 & 0.211796991 & 19384  & 0.034265310& 0.002632326 \\
				4      & 1.846643308 & 0.158557773 & 22384  & 0.031893858& 0.002511350
                \\
				8      & 1.417833461 & 0.114162092 & 25384  & 0.140498287& 0.004407382 \\
				16     & 1.059220795 & 0.080924424 & 32768  & 0.055818578& 0.003344486 \\
				32     & 0.777429162 & 0.057079029 & 40168  & 0.039996452& 0.002554959 \\
				64     & 0.563964349 & 0.040386101 & 47568  & 0.032653545& 0.002321277 \\
				128    & 0.405901600 & 0.028594952 & 54968  & 0.077168816& 0.003802871 \\
				256    & 0.290573182 & 0.020198373 & 65536  & 0.042298414& 0.002611052 \\
				512    & 0.207243911 & 0.014301147 & 85536  & 0.026961063& 0.002003061 \\
				1024   & 0.147431924 & 0.010144281 & 105536 & 0.062300434& 0.002748798 \\
				2048   & 0.104694096 & 0.007218810 & 125536 & 0.031792132& 0.002210619 \\
				4096   & 0.074251823 & 0.005129138 & 131072 & 0.028769748& 0.001954450 \\
				8192   & 0.052614898 & 0.003643280 & 262144 & 0.012148304& 0.001521695 \\
				16384  & 0.037259807 & 0.002747920  & 524288 & 0.007406913& 0.001393969
                \\
				\hline
			\end{tabular}
			\caption{The comparison of theoretical upper bounds and simulated distances for different values of $\lambda$. The simulated distance is computed from 1,000,000 repeated simulations.}
            \label{tab:AI}
		\end{table}

        \newpage
        \begin{comment}
            
        
		The following table presents the bounds obtained from simulations.
		\begin{table}[htbp]
			\centering
			\begin{tabular}{|c|c|c|c|}
				\hline
				$\lambda$ & Simulation upper bounds & $\lambda$ & Simulation upper bounds \\
				\hline
				2      & 0.211796991 & 16384  & 0.002747920 \\
				4      & 0.158557773 & 19384  & 0.002632326 \\
				8      & 0.114162092 & 22384  & 0.002511350 \\
				16     & 0.080924424 & 25384  & 0.004407382 \\
				32     & 0.057079029 & 32768  & 0.003344486 \\
				64     & 0.040386101 & 40168  & 0.002554959 \\
				128    & 0.028594952 & 47568  & 0.002321277 \\
				256    & 0.020198373 & 54968  & 0.003802871 \\
				512    & 0.014301147 & 65536  & 0.002611052 \\
				1024   & 0.010144281 & 85536  & 0.002003061 \\
				2048   & 0.007218810 & 105536 & 0.002748798 \\
				4096   & 0.005129138 & 125536 & 0.002210619 \\
				8192   & 0.003643280 & 131072 & 0.001954450 \\
				16384  & 0.002747920 & 262144 & 0.001521695 \\
				19384  & 0.002632326 & 524288 & 0.001393969 \\
				\hline
			\end{tabular}
			\caption{Simulation results for different values of $\lambda$. For each value of $\lambda$, we generate 1,000,000 samples and repeat this 100 times, then average the results.}
		\end{table}
		\end{comment}
        
		%The following plot shows the simulation results. The graph exhibits three jumps as it exceeds the threshold values. A zoomed-in view around these thresholds is shown in the next figure.
        
Figure \ref{fig:AI} shows the simulation results. The graph exhibits three jumps around the threshold values $n_1,n_2,$ and $n_3$, with zoomed-in views shown in Figure \ref{fig:AI2}. These jumps are also reflected in the bound in  \eqref{sim:AI1}. When $\lambda$ is close to one of the thresholds, $N$ is more likely to exceed the threshold, which results in more effects from the second to the fourth terms in \eqref{sim:AI1}. Moreover, since $r_1>r_2>r_3>r_4,$
we have $\frac{1}{r_1}<\frac{1}{r_2}<\frac{1}{r_3}<\frac{1}{r_4},$
so the mean waiting time increases across the groups, making the jumps more pronounced.


        

		\begin{figure}[H]  
			\centering
            \includegraphics[width=0.8\textwidth]{Fig4.png}
            \caption{Simulated distances as functions of $\lambda$, plotted on a logarithmic scale with consistent $x$-axis spacing.}
            \label{fig:AI}
		\end{figure}


		
\begin{figure}[H]  
			\centering
            \includegraphics[width=0.8\textwidth]{Fig5.png}
            \caption{Zoomed-in view around selected thresholds $n_1 = 25{,}000$, $n_2 = 50{,}000$, and $n_3 = 100{,}000$. The dashed vertical lines indicate these points, where noticeable changes
in the simulated distance can be observed. The figure is based on simulation results.}
        \label{fig:AI2}
		\end{figure}


        
Figure \ref{fig:AI3} compares the simulated values with the theoretical upper bounds. The jumps in the theoretical upper bound are also clearly visible in the plot. Therefore, the simulated results are consistent with the theoretical upper bounds.

		\begin{figure}[H]  % 'h' means here, but you can use 't' for top, 'b' for bottom, etc.
			\centering
            \includegraphics[width=0.8\textwidth]{Fig6.png}
            \caption{Simulated and theoretical upper bounds as functions of $\lambda$, plotted on a logarithmic scale with a consistent $x$-axis spacing.}
            \label{fig:AI3}
		\end{figure}



\section{Conclusion}
\label{sec:sum}

The main contribution of our work is a new upper bound for the $L^1$ distance between a random sum and the standard normal. Although there are comparable results in the literature, our approach is new, and the assumption is different. We apply Stein's method through an approximate zero bias transformation to obtain the bound of the rate \ncolor{$1/\sqrt{n}$} which is the same as the existing results but our constant is larger. However, we drop the identical assumption of the summands. This relaxation is beneficial in real applications since this assumption is often violated. 

We apply our main results to three different real-world cases, namely, simple random sampling with outliers, auto insurance claims, and generative AI wait times. In the latter two cases, we consider settings in which the summands $X_i$ are not identically distributed. Although the actual distance and the corresponding bound are not close when the mean number of summands is small, they tend to be closely aligned as the mean becomes larger. This discrepancy is primarily due to the large constant in the bound, which represents an important direction for future research. These examples demonstrate that relaxing the identical-distribution assumption allows our results to be applied to a broader range of real-world problems, where identical summands are uncommon. Finally, our proposed bounds can be similarly applied to other real-world problems in which the summands are not necessarily identically distributed.

\backmatter

\bmhead{Acknowledgements}

Both authors would like to thank Chutiphan Charoensuk, a research assistant supported by the Grants for Sci Seeding, Faculty of Science, Chulalongkorn University, for assistance with data preparation and numerical simulations. Nathakhun would like to also thank Dawud Thongtha for a fruitful conversation during the initial stage of this project.

%Acknowledgements are not compulsory. Where included they should be brief. Grant or contribution numbers may be acknowledged.

%Please refer to Journal-level guidance for any specific requirements.

\begin{comment}
\bigskip\noindent


\bigskip\noindent
{\bf Data Availability} The artificial datasets used in this study can be obtained from the corresponding author upon request. The PEA dataset is confidential.

\bigskip\noindent
{\bf Ethics approval} The project did not involve ant studies with human or animal participants.

\bigskip\noindent
{\bf Conflict of interest} The authors declare no competing interest.
 

%%===================================================%%
%% For presentation purpose, we have included        %%
%% \bigskip command. Please ignore this.             %%
%%===================================================%%
\begin{comment}
    

\bigskip
\begin{flushleft}%
Editorial Policies for:

\bigskip\noindent
Springer journals and proceedings: \url{https://www.springer.com/gp/editorial-policies}

\bigskip\noindent
Nature Portfolio journals: \url{https://www.nature.com/nature-research/editorial-policies}

\bigskip\noindent
\textit{Scientific Reports}: \url{https://www.nature.com/srep/journal-policies/editorial-policies}

\bigskip\noindent
BMC journals: \url{https://www.biomedcentral.com/getpublished/editorial-policies}
\end{flushleft}
\end{comment}


\section*{Declarations}

\noindent\textbf{Declaration of generative AI in the writing process:}
During the preparation of this manuscript, the authors used ChatGPT (OpenAI) solely for language editing, grammar checking, and improving the clarity of the manuscript. The authors reviewed and approved all AI-assisted text and take full responsibility for the final content of the manuscript.

\noindent\textbf{Funding:} Not applicable.

\noindent\textbf{Conflict of interest/Competing interests:} The author declares no conflict of interest.

\noindent\textbf{Data availability:} Data are available upon request.




%\begin{appendices}


%\end{appendices}

%%===========================================================================================%%
%% If you are submitting to one of the Nature Portfolio journals, using the eJP submission   %%
%% system, please include the references within the manuscript file itself. You may do this  %%
%% by copying the reference list from your .bbl file, paste it into the main manuscript .tex %%
%% file, and delete the associated \verb+\bibliography+ commands.                            %%
%%===========================================================================================%%

\bibliography{sn-bibliography}% common bib file
%% if required, the content of .bbl file can be included here once bbl is generated
%%\input sn-article.bbl

\end{document}
